\documentclass[a4paper,11pt,bibliography=totoc,listof=totoc,headinclude=true,open=right]{scrbook}
% Option "oneside" für einseitigen Druck. Weglassen, falls die Arbeit doppelseitig gedruckt wird

\usepackage[english, ngerman]{babel}
\usepackage[utf8]{inputenc}
%\usepackage{fullpage}
\usepackage{ifthen}
\usepackage{color}
\usepackage{amsmath,amsthm,amssymb,amsfonts,stmaryrd}
\usepackage{graphicx}
\usepackage{psfrag}
\usepackage{setspace}
\usepackage[shortlabels]{enumitem}
\usepackage{mathtools}
\usepackage{bbm}
\usepackage{mathrsfs}
\usepackage{tikz}

% links in pdf
\usepackage[unicode,colorlinks=false,pagebackref=false]{hyperref}

\usepackage{cleveref}

% Zum Druck verwende schwarze Links!
%\usepackage[unicode,colorlinks=true,linkcolor=black,citecolor=black,urlcolor=black,pagebackref=false]{hyperref} 
	% colorlinks=false umrahmt Links statt einzufaerben, 

\nocite{*}

% document style
\KOMAoptions{footinclude=false} % Fusszeile wird nicht zu Satzspiegel gezaehlt
\KOMAoptions{headsepline=true} % Trennlinie zwischen Kopfzeile und Text
\KOMAoptions{DIV=12} % beeinflusst Satzspiegel
\KOMAoptions{BCOR=8mm} % Bindekorrektur
\pagestyle{headings} % mit Kopfzeilen
\addtokomafont{disposition}{\rmfamily}
\setlength{\parindent}{0pt}
\setlength{\parskip}{6pt}

\raggedbottom

\newcommand{\todo}[1]{{\color{todo} #1}}
\renewcommand{\qedsymbol}{$\blacksquare$}
% \newcommand{\cl}[1]{\operatorname{cl}(#1)}
\newcommand{\cl}[1]{c\ell({#1})}
\newcommand{\N}{\mathbb{N}}
\newcommand{\Q}{\mathbb{Q}}
\newcommand{\R}{\mathbb{R}}
\newcommand{\C}{\mathbb{C}}
\newcommand{\Z}{\mathbb{Z}}
\newcommand{\D}{\mathbb{D}}
\newcommand{\F}{\mathscr{F}}
\renewcommand{\O}{\mathcal{O}}
\newcommand{\powerseries}[2][n]{\sum_{#1=0}^\infty {#2}z^{#1}}
\newcommand{\diff}{\mathop{}\!\mathrm{d}}
\newcommand{\deriv}[2]{\frac{\mathrm{d}^{#2}}{\mathrm{d} {#1}^{#2}}}
\newcommand{\1}{\mathbbm{1}}
\newcommand{\restr}[2]{{% we make the whole thing an ordinary symbol
  \left.\kern-\nulldelimiterspace % automatically resize the bar with \right
  #1 % the function
  \vphantom{\big|} % pretend it's a little taller at normal size
  \right|_{#2} % this is the delimiter
}}
\newcommand{\mapping}[3]{
  \ensuremath{
    #1 : \left\{
    \begin{array}{@{\;}c@{\;}c@{\;}l@{}}
      #2 \\
      #3
    \end{array}
    \right.
  }
}

\renewcommand{\Re}{\operatorname{Re}}
\renewcommand{\Im}{\operatorname{Im}}

\theoremstyle{plain}
\newtheorem{theorem}{Theorem}[section]
\newtheorem{lemma}[theorem]{Lemma}
\newtheorem{corollary}[theorem]{Corollary}
\newtheorem{proposition}[theorem]{Proposition}

\theoremstyle{definition}
\newtheorem{definition}[theorem]{Definition}
\newtheorem{remark}[theorem]{Remark}
\newtheorem{example}[theorem]{Example}

\AtBeginEnvironment{remark}{%
  \pushQED{\qed}\renewcommand{\qedsymbol}{$\sslash$}
}
\AtEndEnvironment{remark}{\popQED\endremark}

\AtBeginEnvironment{definition}{%
  \pushQED{\qed}\renewcommand{\qedsymbol}{$\sslash$}
}
\AtEndEnvironment{definition}{\popQED\enddefinition}

\AtBeginEnvironment{example}{%
  \pushQED{\qed}\renewcommand{\qedsymbol}{$\sslash$}
}
\AtEndEnvironment{example}{\popQED\endexample}

\newlist{arrowlist}{itemize}{1}
\setlist[arrowlist]{label=$\rightsquigarrow$}

\recalctypearea % berechne Satzspiegel neu

\definecolor{change}{rgb}{0,.55,.55}
\definecolor{todo}{rgb}{.75,0.2,0.2}

\def\revision#1{{\color{red}#1}}

%%%%%%%%%%%%%%%%%%%%%%%%%%%%%%%%%%%%%%%%%%%%%%%%%%%%%%%%%%%%%%%%%%%%%%%%%%%%%%%%%%%%%%%%%%%%%%%%%%%%%%%%%%%%%%
%%%%%%%%%%%%%%%%%%%%%%%%%%%%%%%%%%%%%%%%%%%%%%%%%%%%%%%%%%%%%%%%%%%%%%%%%%%%%%%%%%%%%%%%%%%%%%%%%%%%%%%%%%%%%%
%%%%%%%%%%%%%%%%%%%%%%%%%%%%%%%%%%%%%%%%%%%%%%%%%%%%%%%%%%%%%%%%%%%%%%%%%%%%%%%%%%%%%%%%%%%%%%%%%%%%%%%%%%%%%%
%%%%%%%%%%%%%%%%%%%%%%%%%%%%%%%%%%%%%%%%%%%%%%%%%%%%%%%%%%%%%%%%%%%%%%%%%%%%%%%%%%%%%%%%%%%%%%%%%%%%%%%%%%%%%%

\begin{document}

%%%%%%%%%%%%%%%%%%%%%%%%%%%%%%%%%%%%%%%%%%%%%%%%%%%%%%%%%%%%%%%%%%%%%%%%%%%%%%%%%%%%%%%%%%%%%%%%%%%%%%%%%%%%%%
% TITELSEITE [OBLIGATORISCH]
%%%%%%%%%%%%%%%%%%%%%%%%%%%%%%%%%%%%%%%%%%%%%%%%%%%%%%%%%%%%%%%%%%%%%%%%%%%%%%%%%%%%%%%%%%%%%%%%%%%%%%%%%%%%%%

\pagenumbering{Alph}
\selectlanguage{ngerman}

\begin{titlepage}
  %\vspace*{-2cm}
  \begin{center}
    \includegraphics[width=0.45\textwidth]{TULogo.eps}
    \vskip 1cm%
    {\LARGE B~\Large A~C~H~E~L~O~R~A~R~B~E~I~T}
    \vskip 8mm
    {\huge\bfseries\setstretch{1.1} Picard's Great Theorem \\ \vspace{-0.1cm}{\Large and} \\ Growth, Zeros, and Composition \\ of Entire Functions \par}
    \vskip 1cm
    \large 
    ausgef\"uhrt am    
    \vskip 0.75cm
    {\Large Institut f\"ur\\[1ex] Analysis und Scientific Computing}\\[1ex]
    {\Large TU Wien}
    \vskip0.75cm
    unter der Anleitung von
    \vskip0.75cm
    {\Large\bfseries Ao.Univ.Prof. Dipl.-Ing. Dr.techn. Harald Woracek}\\[1ex]
    \vskip 0.5cm
    durch
    \vskip 0.5cm
    {\Large\bfseries Ian Hornik}\\[1ex]
    {Matrikelnummer: 12115564}\\[1ex]
    % {Haizingergasse 18/4}\\[1ex]
    % {1180 Wien}
  \end{center}
  
  \vfill
  
  \small
  Wien, am \today
  \vspace*{-15mm}
\end{titlepage}

\cleardoublepage

\pagenumbering{roman}
\selectlanguage{english} 

\tableofcontents

\cleardoublepage
\pagenumbering{arabic}

\chapter{Introduction}
\label{ch:introduction}

In this thesis, we explore the theory of entire functions, that is, complex-valued functions that are holomorphic in the entire complex plane.

Our early focus lies on Picard's Great Theorem, an important corollary of which states that an entire function assumes every complex value infinitely often, with at most one exception. A natural example for this behaviour is the complex exponential function, which does not assume zero, but does assume all other values infinitely often. \Cref{ch:picards-great-theorem} is dedicated to deriving this theorem and providing groundwork for the results in the later chapters.

\Cref{sec:order,sec:type} then briefly introduce a scale of asymptotic growth for the modulus of entire functions, in the form of order and type. Roughly speaking, if $f$ is of finite order, then $f$ grows no faster than a function of the form $e^{Q(z)}$, where $Q$ is a polynomial, as $\vert z \vert \to \infty$. Results introduced here are not particularly interesting, but will frequently be utilized towards the end of the chapter.

\Cref{sec:zeros,sec:canonical-products,sec:hadamards-theorem} explore the nature of zeros of entire functions. As is well-known, a complex polynomial can be represented solely by its zeros and some scaling factor. In \Cref{sec:canonical-products}, this property is generalized for entire functions by the Weierstraß Factorization Theorem: For any entire function $f$ there exists an entire function $g$ and a so-called ``canonical product'' $\Pi$, which is determined by the zeros of $f$, such that
$$ f(z) = z^m e^{g(z)} \Pi(z), \quad \textrm{for all } z \in \C, $$
where $m \in \N_0$ denotes the order of the zero of $f$ at the origin. Hadamard's Theorem, which is the focus of \Cref{sec:hadamards-theorem}, then provides an immediate refinement for entire functions of finite order $\rho $, showing that $g$ can be taken to be a polynomial of degree at most $\rho$. This has immediate corollaries, in particular if the order is not an integer, which we will explore towards the end of the section.

\todo{I should clarify what ``degree at most'' means if $\rho$ is not an integer. Maybe note which values $\rho$ can take overall.}

Finally, \Cref{ch:composition-entire-functions} studies when the composition of functions of finite order is, again, of finite order. Pólya's Theorem will provide necessary conditions for this to hold. A consequence is Thron's Theorem, which states that for any entire function $g$ of finite order, which assumes some complex value only finitely often, there does not exist an entire function $f$ for which $f \circ f = g$ holds on the entire complex plane.

Introductory results on complex analysis are assumed as prerequisites and will be used as found in \cite{bluemlinger-complex-analysis,stein-shakarchi-princeton}. These include Cauchy's integral formula, the maximum modulus principle, the theorems of Montel, Hurwitz and Jensen, the Schwarz lemma and standard results on power series.

The following is an overview of the notation used:

\begin{itemize}[label=$\rightsquigarrow$]
    \item The sets $\N$ and $\Z$ denote, respectively, the natural numbers (starting with $1$) and the integers. Furthermore we define $\N_0 \coloneqq \N \cup \{ 0 \}$.
    \item For a set $A$, its indicator function $\1_A(a)$ is defined as $1$ if $a \in A$ and $0$ otherwise.
    \item The restriction of a function $f : M \to N$ to a set $D \subseteq M$ will be denoted $\restr{f}{D}$.
    \item The sets $\R$ and $\C$ denote, respectively, the real line and the complex plane, which will both be endowed with the standard Euclidean metric and topology.
    \item The closure of a subset $A$ of $\R$ or $\C$ (with respect to the appropriate topology) will be denoted as $\cl{A}$ and the boundary as $\partial A$.
    \item For subsets $M, N$ of $\R$ or $\C$ we denote by $C(M, N)$ the space of continuous maps from $M$ into $N$.
    \item For a real- or complex-valued function $f$ defined on some set $M$ we define $\Vert f \Vert_M \coloneqq \sup_{x \in M} \vert f(x) \vert$.
    \item If a function $f$ is differentiable, in the real or complex sense, then $f'$ and $f''$ denote the respective first and second derivatives of $f$. For $n \in \N$, the $n$-th derivative will be denoted $f^{(n)}$.
    \item Let $a \in \R$ and let $f, g$ be functions on $[a, \infty)$ such that $f$ is either real- or complex-valued and $g$ is real-valued. Then we write $f(x) = \O(g(x))$ as $x \to \infty$ if there exist constants $C > 0$ and $x_0 \geq a$ such that for $x \geq x_0$ we have $\vert f(x) \vert \leq C g(x)$.
    \item For $a, b \in \C$ we denote by $[a, b]$ the line segment connecting $a$ and $b$, parametrized by $\gamma : [0, 1] \to \C, t \mapsto (1 - t) a + tb$.
    \item If $G \subseteq \C$, then $G$ is said to be a \emph{domain} if $G$ is non-empty, open and connected.
    \item For a point $a \in \C$ and $r > 0$ we denote by $B_r(a)$ the open disk of radius $r$ centered at $a$. The open unit disk $B_1(0)$ will be denoted $\D$. Furthermore we define $\C^\times \coloneqq \C \setminus \{ 0 \}$ and $\D^\times \coloneqq \D \setminus \{ 0 \}$.
    \item Let $G \subseteq \C$ be open, then a function $f : G \to \C$ is said to be \emph{holomorphic} at $z \in G$ if $f$ is complex differentiable at $z$. If $f$ is holomorphic at all $z \in G$, then $f$ is said to be holomorphic on $G$. The space of all such holomorphic functions on $G$ will be denoted as $H(G)$. If $f \in H(\C)$, then $f$ is called \emph{entire}.
    \item If $G \subseteq \C$ is a domain and $D \subseteq \C$ is a set with $\cl{D} \subseteq G$, then the contour integral of a function $f \in H(G)$ along the boundary of $D$ will be denoted
    $$ \oint_{\partial D} f(z) \diff z $$
    and is, unless otherwise indicated, always parametrized with mathematically positive (counter-clockwise) orientation. \todo{Only for NICE boundaries.}
    \item The exponential function will be denoted $\exp : \C \to \C^\times$. The (local) inverse will be denoted $\log$, where we, unless specified otherwise, always choose the principle branch, i.e., the branch where $\log 1 = 0$.
\end{itemize}

\chapter{Picard's Great Theorem}
\label{ch:picards-great-theorem}

\todo{Introduction, mention E. Picard, essential singularities, Casorati-Weierstrass.}

\section{Bloch's Theorem}
\label{sec:blochs-theorem}

\begin{lemma} \label{lem:bloch-lemma-1}
    Let $G \subset \C$ be a bounded domain, $f : \cl{G} \to \C$ continuous and $\restr{f}{G} : G \to \C$ open. If there exists $a \in G$ such that $s \coloneqq \min_{z \in \partial G} \vert f(z) - f(a) \vert > 0$, then $B_{s}(f(a)) \subseteq f(G)$.
\end{lemma}

\begin{proof}
    The function $z \mapsto \vert z - f(a) \vert$ is continuous on the compact set $\partial f(G)$, hence it attains its minimum $m$ at some $w_* \in \partial f(G)$. Let $(z_n)_{n \in \N}$ be a sequence in $G$, such that $\lim_{n \to \infty} f(z_n) = w_*$. Since $\cl{G}$ is compact, we can find a subsequence that converges to some $z_* \in \cl{G}$ and by continuity of $f$ we have $f(z_*) = w_*$.

    Assuming $z_* \in G$, since $\restr{f}{G}$ is open, the image of any open set in $G$ containing $z_*$ under $f$ must be an open set in $f(G)$ containing $w_*$, which is impossible since $w_* \in \partial G$.

    Therefore $z_* \in \partial G$ and we have $m = \vert w_* - f(a) \vert = \vert f(z_*) - f(a) \vert \geq s$, from which it follows that $B_{s}(f(a)) \subseteq f(G)$.
\end{proof}

\begin{lemma} \label{lem:bloch-lemma-2}
    Let $a \in \C, r > 0$ and $B \coloneqq B_{r}(a)$. Suppose further that $f \in H(\cl{B})$ such that $\Vert f' \Vert_B \leq 2 \vert f'(a) \vert$. Then $ B_{R}(f(a)) \subseteq f(B)$, where $ R \coloneqq (3 - 2 \sqrt{2}) r \vert f'(a) \vert $.
\end{lemma}

\begin{proof}
    Without loss of generality we can assume $a = f(a) = 0$. Consider the function
    $$ \mapping{\alpha_f}{B & \to & \C,}{z & \mapsto & f(z) - f'(0) z,} $$
    which satifies, for the line $\gamma$ connecting $0$ to $z \in B$,
    \begin{equation} \label{eq:bloch-lemma-2:estimate-1}
        \vert \alpha_f(z) \vert = \left\vert \int_\gamma f'(\zeta) - f'(0) \diff \zeta \right\vert \leq \int_0^1 \vert f'(t z) - f'(0) \vert \vert z \vert \diff t. \tag{\textasteriskcentered}
    \end{equation}
    Let $w \in B$, then Cauchy's integral formula gives
    \begin{align*}
        \vert f'(w) - f'(0) \vert &= \frac{1}{2\pi} \left\vert \int_{\partial B} \frac{f'(\zeta)}{\zeta - w} - \frac{f'(\zeta)}{\zeta} \diff \zeta \right\vert = \frac{1}{2\pi} \left\vert \int_{\partial B} \frac{w f'(\zeta)}{\zeta (\zeta - w)} \diff \zeta \right\vert \leq \\
        &\leq \frac{1}{2\pi} \int_{\partial B} \frac{\vert w \vert \Vert f' \Vert_B}{r(r - \vert w \vert)} \diff \zeta = \frac{\vert w \vert}{r - \vert w \vert} \Vert f' \Vert_B.
    \end{align*}
    Combining the above with \eqref{eq:bloch-lemma-2:estimate-1} and our estimate on $\Vert f' \Vert_B$ yields
    \begin{align*}
        \vert \alpha_f(z) \vert &\leq \int_0^1 \frac{\vert z t \vert \Vert f' \Vert_B}{r - \vert z t \vert} \vert z \vert \diff t \leq \frac{\vert z \vert^2}{r - \vert z \vert} \Vert f' \Vert_B \int_0^1 t \diff t \leq \frac{\vert z \vert^2}{r - \vert z \vert} \vert f'(0) \vert.
    \end{align*}
    Let $0 < \rho < r$, then for $\vert z \vert = \rho$ we have
    \begin{align*}
        \vert f'(0) \vert \rho - \vert f(z) \vert &\leq \vert \alpha_f(z) \vert \leq \frac{\rho^2}{r - \rho} \vert f'(0) \vert \\
        \Longleftrightarrow \quad \vert f(z) \vert &\geq \left( \rho - \frac{\rho^2}{r - \rho} \right) \vert f'(0) \vert.
    \end{align*}
    Considering $\rho - \rho^2 / (r - \rho)$ as a function of $\rho$, it attains its maximum value, $(3 - 2 \sqrt{2}) r$, at $\rho_* \coloneqq (1 - \sqrt{2} / 2)r \in (0, r)$. Therefore,
    \begin{equation*}
        \vert f(z) \vert \geq (3 - 2 \sqrt{2}) r \vert f'(0) \vert = R, \quad \textrm{for all} \quad \vert z \vert = \rho_*.
    \end{equation*}
    Invoking \Cref{lem:bloch-lemma-1} with $G \coloneqq B_{\rho_*}(0)$ thus yields $B_{R}(0) \subseteq f(G) \subseteq f(B)$.
\end{proof}

\begin{theorem} \label{thm:bloch-stronger}
    Let $f \in H(\cl{\D})$ be non-constant. Then there is a point $p \in \D$ and a constant $C_f > 0$ such that $B_{R}(f(p)) \subseteq f(\D)$, where $R \coloneqq (\frac{3}{2} - \sqrt{2}) C_f \geq (\frac{3}{2} - \sqrt{2}) \vert f'(0) \vert$.
\end{theorem}

\begin{proof}
    The function
    $$ \mapping{\alpha_f}{\cl{\D} & \to & \R}{z & \mapsto & \vert f'(z) \vert (1 - \vert z \vert)} $$
    is continuous and attains its maximum $C_f > 0$ at some point $p \in \D$.
    
    Set $t \coloneqq \frac{1}{2}(1 - \vert p \vert) > 0$, then we have $B_{t}(p) \subseteq \D$ and $1 - \vert z \vert \geq t$ for $z \in B_{t}(p)$. Since $\vert f'(z) \vert (1 - \vert z \vert) \leq C_f = 2 t \vert f'(p) \vert$, this implies $\vert f'(z) \vert \leq 2 \vert f'(p) \vert$. By \Cref{lem:bloch-lemma-2}, we obtain
    $ B_{R}(f(p)) \subseteq f(\D) $, where 
    $$ R \coloneqq (3 - 2 \sqrt{2}) t \vert f'(p) \vert = ({\textstyle \frac{3}{2}} - \sqrt{2}) \vert f'(p) \vert (1 - \vert p \vert) > {\textstyle \frac{1}{12}} \vert f'(0) \vert, $$
    establishes the assertion.
\end{proof}

We immediately have:    

\begin{theorem}[Bloch] \label{thm:bloch}
    Let $f \in H(\cl{\D})$ and assume that $f'(0) = 1$. Then $f(\D)$ contains a disk of radius $\frac{3}{2} - \sqrt{2}$.
\end{theorem}

In the following we will denote by $\beta > 0$ any constant satisfying Bloch's Theorem, for example $\beta = \frac{1}{12} < \frac{3}{2} - \sqrt{2}$.

\begin{corollary} \label{cor:bloch-domain}
    Let $G \subset \C$ be a domain and $f \in H(G)$ with $f'(c) \neq 0$ for some $c \in G$. Then $f(G)$ contains a disk of every radius $\beta s \vert f'(c) \vert$, where $0 < s < d(c, \partial G)$.
\end{corollary}

\begin{proof}
    Without loss of generality we may assume $c = 0$. Since $f$ is analytic on $\cl{B_s(0)} \subseteq G$, we have $g(z) \coloneqq f(sz) / sf'(0) \in H(\cl{\D})$. Since $g'(0) = 1$, Bloch's Theorem (\ref{thm:bloch}) yields a disk $B$ of radius $\beta$ such that $B \subseteq g(\D)$. It follows that
    \begin{equation*}
        s \vert f'(0) \vert B \subseteq s \vert f'(0) \vert g(\D) = f(B_s(0)) \subseteq f(G).
    \end{equation*}
\end{proof}

\begin{corollary} \label{cor:bloch-entire}
    If $f \in H(\C)$ is non-constant, then $f(\C)$ contains a disk of every radius.
\end{corollary}
\section{Schottky's Theorem}
\label{sec:schottkys-theorem}

Holomorphic functions which omit the values $0$ and $1$ have a universal estimate on the growth of their modulus, which will be given by Schottky's Theorem.

For a domain $G \subseteq \C$ and a set $E \subseteq \C$ we define $H(G; E)$ as the set\footnote{Note that, unlike $H(G)$, the set $H(G; E)$ is usually not a linear space.} of all $f \in H(G)$ such that $f(G) \subseteq E$.

\begin{lemma} \label{lem:schottky-1}
    It holds that:
    \begin{enumerate}[i.]
        \item If $a, b \in \R$ with $\cos \pi a = \cos \pi b$, then $b = \pm a + 2n$ for some $n \in \Z$.
        \item For every $w \in \C$ there exists $v \in \C$ such that $\cos \pi v = w$ and $\vert v \vert \leq 1 + \vert w \vert$.
    \end{enumerate}
\end{lemma}

\begin{proof}
    For the first part it suffices to note that
    $$ 0 = \cos \pi a - \cos \pi b = \textstyle -2 \sin \frac{\pi}{2} ( a + b ) \sin \frac{\pi}{2} ( a - b ). $$
    Since the complex cosine function is surjective and $2 \pi$-periodic, we can choose $v = a + i b$ with $w = \cos \pi v$ and $\vert a \vert \leq 1$. Now we have
    \begin{align*}
        \vert w \vert^2 &= \vert \cos (\pi a + i \pi b) \vert^2 = \vert \cos \pi a \cos i \pi b + \sin \pi a \sin i \pi b \vert^2 = \\
        &= \vert \cos \pi a \cosh \pi b - i \sin \pi a \sinh \pi b \vert^2 = \\
        &= \cos^2 \pi a \cosh^2 \pi b + \sin^2 \pi a \sinh^2 \pi b = \\
        &= \cos^2 \pi a + \cos^2 \pi a \sinh^2 \pi b + \sin^2 \pi a \sinh^2 \pi b = \\
        &= \cos^2 \pi a + \sinh^2 \pi b \geq \sinh^2 \pi b \geq \pi^2 b^2,
    \end{align*}
    where the last inequality holds since $\sinh^2 x \geq x^2$ for $x \in \R$. We conclude
    \begin{equation*}
        \vert v \vert = \sqrt{a^2 + b^2} \leq \sqrt{1 + \vert w \vert^2 / \pi^2} \leq 1 + \vert w \vert.
    \end{equation*}
\end{proof}

We recall the following result: Let $G \subseteq \C$ be a simply connected domain and $f \in H(G)$, such that $f$ vanishes nowhere on $G$. There is a function\footnote{$g$ is also sometimes called a \emph{logarithm} of $f$.} $g \in H(G)$ such that $f = e^g$, which, for $n \in \N$, can also be used to obtain $n$-th roots of such functions by defining $\sqrt[n]{f} \coloneqq e^{g/n}$.

\begin{lemma} \label{lem:schottky-2}
    Let $G \subseteq \C$ be a simply connected domain and $f \in H(G; \C \setminus \{ -1, 1 \})$, then there exists $F \in H(G)$ such that $f = \cos F$.
\end{lemma}

\begin{proof}
    Since $1 - f^2$ vanishes nowhere in $G$, it has a square root $g \in H(G)$, therefore
    \begin{equation*}
        1 = f^2 + g^2 = (f + ig)(f - ig).
    \end{equation*}
    Thus $f + ig$ vanishes nowhere and there exists an $F \in H(G)$ with $f + ig = e^{iF}$. By the above we also have $f - ig = e^{-iF}$ and therefore
    \begin{equation*}
        f = {\textstyle \frac{1}{2}} (e^{iF} + e^{-iF}) = \cos F. \qedhere
    \end{equation*}
\end{proof}

\begin{lemma} \label{lem:schottky-3}
    Let $G \subseteq \C$ be a simply connected domain and $f \in H(G; \C \setminus \{ 0, 1 \})$. Then there exists $g \in H(G)$ such that:
    \begin{enumerate}[i.]
        \item $f = \frac{1}{2} ( 1 + \cos \pi (\cos \pi g))$.
        \item $\vert g(0) \vert \leq 3 + 2 \vert f(0) \vert$.
        \item $g(G)$ contains no disk of radius $1$.
        \item If $\D \subseteq G$, then $\vert g(z) \vert \leq \vert g(0) \vert + \frac{\theta}{\beta (1 - \theta)}$ for all $\vert z \vert \leq \theta$, where $0 < \theta < 1$.
    \end{enumerate}
\end{lemma}

\begin{proof}
    By \Cref{lem:schottky-2}, there exists $\widetilde{F} \in H(G)$ such that $2f - 1 = \cos \pi \widetilde{F}$ and, by \Cref{lem:schottky-1}, there is $b \in \C$ with $\cos \pi b = 2f(0) - 1$ and $\vert b \vert \leq 1 + \vert 2 f(0) - 1 \vert \leq 2 + 2 \vert f(0) \vert$. Furthermore, since $\cos \pi b = \cos \pi \widetilde{F}(0)$, we have $b = \pm \widetilde{F}(0) + 2k$ for some $k \in \Z$. Then $F \coloneqq \pm \widetilde{F} + 2k \in H(G)$ satisfies $F(0) = b$ and $2f - 1 = \cos \pi F$.
    
    Since $F$ must omit all integers, there exists $\widetilde{g} \in H(G)$ such that $F = \cos \pi \widetilde{g}$. Similarly, there is $a \in \C$ such that $\cos \pi a = b$ and $\vert a \vert \leq 1 + \vert b \vert \leq 3 + 2 \vert f(0) \vert$. Like in the above, since $\cos \pi a = \cos \pi \widetilde{g}(0)$, we have $a = \pm \widetilde{g}(0) + 2\ell$ for some $\ell \in \Z$, thus $g \coloneqq \pm \widetilde{g} + 2\ell \in H(G)$ satisfies $g(0) = a$ and $F = \cos \pi g$. Together we obtain
    $$ \textstyle f = \frac{1}{2} (1 + \cos \pi (\cos \pi g)), \quad \textrm{and} \quad \vert g(0) \vert = \vert a \vert \leq 3 + 2 \vert f(0) \vert $$
    and thus have shown i. and ii.

    To show iii. we consider the set
    $$ A \coloneqq \{ m \pm i \pi^{-1} \log (n + \sqrt{n^2 - 1}) : m \in \Z, n \in \N \setminus \{ 0 \} \}, $$
    the points of which can be considered the vertices of a rectangular grid in $\C$, the cells of which are of fixed width and variable height. The width of such a rectangular cell is $1$, and since
    \begin{align*}
        \log ((n+1) &+ \sqrt{(n+1)^2 - 1}) - \log (n + \sqrt{n^2 - 1}) = \\
        &= \log \frac{1 + \frac{1}{n} + \sqrt{1 + \frac{2}{n}}}{1 + \sqrt{1 - \frac{1}{n^2}}} \leq \log (2 + \sqrt{3}) < \pi,
    \end{align*}
    their height is bounded above by some $C < 1$. Therefore, for all $z \in \C$ there is $w_z \in A$ such that $\vert \Re z - \Re w_z \vert \leq \frac{1}{2}$ and $\vert \Im z - \Im w_z \vert \leq \frac{C}{2}$. Thus, we have
    $$ \vert z - w_z \vert \leq \vert \Re z - \Re w_z \vert + \vert \Im z - \Im w_z \vert \leq \frac{1}{2} + \frac{C}{2} < 1. $$
    If we can show that $g(G) \cap A = \emptyset$, then $g(G)$ cannot contain a disk of radius $1$. Let $a = p + i \pi^{-1} \log(q + \sqrt{q^2 - 1}) \in A$, then
    \begin{align*}
        \cos \pi a &= {\textstyle\frac{1}{2}}( e^{i \pi a} + e^{-i \pi a}) = {\textstyle\frac{1}{2}} (-1)^p ((q + \sqrt{q^2 - 1})^{-1} + (q + \sqrt{q^2 - 1})) = \\
        &= \frac{(-1)^p}{2} \frac{1 + q^2 + 2q \sqrt{q^2 - 1} + q^2 - 1}{q + \sqrt{q^2 - 1}} = (-1)^p q
    \end{align*}
    and thus $\cos \pi (\cos \pi a) = \pm 1$. But $0, 1 \notin f(G)$, therefore $a \notin g(G)$ and $g(G) \cap A = \emptyset$, proving (iii).

    For (iv), if $\D \subseteq G$, then $\restr{g}{\D} \in H(\D)$. Fix $0 < \theta < 1$, then for $\vert z \vert \leq \theta$ we have
    $$ d(z, \partial \D) = \inf_{w \in \partial \D} \vert z - w \vert \geq \inf_{w \in \partial \D} \left( \vert w \vert - \vert z \vert \right) \geq 1 - \theta. $$
    From (iii) it follows that $\restr{g}{\D}(\D)$ does not contain a disk of radius $1$. Let $0 < s < 1 - \theta$, then applying \Cref{cor:bloch-domain} to $\restr{g}{\D}$ implies that $\beta s \vert g'(z) \vert < 1$, for any $z \in \D$. Taking the supremum over $s$ and rearranging yields $ \vert g'(z) \vert \leq 1 / (\beta (1 - \theta))$, thus our desired estimate is shown by
    \begin{align*}
        \vert g(z) \vert &\leq \vert g(0) \vert + \vert g(z) - g(0) \vert \leq \vert g(0) \vert + \int_{[0, z]} \vert g'(\zeta) \vert \diff \zeta \leq \vert g(0) \vert + \frac{\theta}{\beta (1 - \theta)}. \qedhere
    \end{align*}
\end{proof}

\begin{theorem}[Schottky] \label{thm:schottky}
    There exists a function ${\psi : (0, 1) \times (0, \infty) \to (0, \infty)}$ such that for any $f \in H(\D; \C \setminus \{ 0, 1 \})$ with $\vert f(0) \vert \leq \omega$ it holds that
    \begin{equation}
        \vert f(z) \vert \leq \psi(\theta, \omega), \quad \vert z \vert \leq \theta.
    \end{equation}
\end{theorem}

\begin{proof}
    Note that for all $w \in \C$ we have $\vert \cos w \vert \leq e^{\vert w \vert}$ and $\frac{1}{2} \vert 1 + \cos w \vert \leq e^{\vert w \vert}$. Hence, from \Cref{lem:schottky-3}, we get
    \begin{align*}
        \vert f(z) \vert &= \vert {\textstyle \frac{1}{2}} ( 1 + \cos \pi ( \cos \pi g(z) ) ) \vert \leq \exp ( \pi \exp \pi \vert g(z) \vert ) \leq \\ &\leq \exp \left(\pi \exp \pi \left( \vert g(0) \vert + \frac{\theta}{\beta(1 - \theta)} \right) \right) \leq \\
        &\leq \exp \left( \pi \exp \pi \left( 3 + 2 \omega + \frac{\theta}{\beta(1 - \theta)} \right) \right),
    \end{align*}
    and defining $\psi(\theta, \omega)$ as the final term establishes the assertion.
\end{proof}

\Cref{lem:schottky-3} is quite powerful, as it contains not only Schottky's Theorem, but Picard's Little Theorem as well:

\begin{theorem}[Picard's Little Theorem] \label{thm:picard-little}
    Let $f \in H(\C; \C \setminus \{ a, b \})$ for distinct points $a, b \in \C$. Then $f$ is constant.
\end{theorem}

\begin{proof}
    Consider $f_1(z) \coloneqq \frac{f(z) - a}{b - a} \in H(\C; \C \setminus \{ 0, 1 \})$. By \Cref{lem:schottky-3} there is some $g \in H(\C)$ such that $f_1 = \frac{1}{2}(1 + \cos \pi (\cos \pi g))$ and $g(\C)$ does not contain a disk of radius $1$. By \Cref{cor:bloch-entire}, we thereby have that $g$ must be constant, and therefore so are $f_1$ and $f$.
\end{proof}
\section{Normal families}
\label{sec:normal-families}

First, we recall a generalized form of locally uniform convergence:

\begin{definition}
    Let $G \subseteq \C$ be a domain, $f \in H(G)$ and $(f_n)_{n \in \N}$ a sequence in $H(G)$. We say that \emph{$f_n$ converges compactly in $G$ to $f$}, or \emph{$f_n$ converges compactly in $G$ to $\infty$} as $n \to \infty$, if for every compact set $K \subset G$
    \begin{equation}
        \lim_{n \to \infty} \sup_{z \in K} \vert f_n(z) - f(z) \vert = 0, \quad \textrm{or} \quad \lim_{n \to \infty} \inf_{z \in K} \vert f_n(z) \vert = \infty. \qedhere
    \end{equation}
\end{definition}

\begin{remark}
    One can define compact convergence generally for functions from a topological space $(X, \mathcal{T})$ into a metric space $(Y, d_Y)$. Compact convergence to $\infty$ can then be seen as compact convergence to a constant function with value $\infty$, where $d_Y(y, \infty)$ is appropriately defined for $y \in Y$. This is precisely the case when considering the chordal metric of the Riemann sphere, but we shall not elaborate further on this.
\end{remark}

\begin{definition}
    Let $f : \C \to \C$ and $a \in \C$, then the \emph{a-points of $f$} are defined as the zeros of $f(z) - a$, that is the set of all points $w \in \C$ with $f(w) = a$.
\end{definition}

A well-known theorem on compact convervence is:

\begin{theorem}[Hurwitz] \label{thm:hurwitz}
    Let $G \subseteq \C$ be a domain and $(f_n)_{n \in \N}$ a sequence in $H(G)$ that converges compactly to $f \in H(G)$. If for every $n \in \N$ the number of $a$-points of $f_n$ is bounded by some $m \in \N_0$, then either
    \begin{itemize}
        \item the number of $a$-points of $f$ are also bounded by $m$ or
        \item $f \equiv a$.
    \end{itemize}
\end{theorem}

As an immediate consequence we obtain that compact convergence is, in some ways, compatible with reciprocals:

\begin{lemma} \label{lem:compact-convergence-reciprocals}
Let $G \subseteq \C$ be a domain and $(f_n)_{n \in \N}$ a sequence in $H(G; \C \setminus \{ 0 \})$. If it converges compactly to some $f \in H(G)$, then it holds that either:
\begin{itemize}
    \item $0 \notin f(G)$ and $1 / f_n \to 1 / f$ compactly in $G$, or
    \item $f \equiv 0$ and $1 / f_n \to \infty$ compactly in $G$.
\end{itemize}
If it converges compactly to $\infty$, then $1 / f_n \to 0$ compactly in $G$.
\end{lemma}

\begin{proof}
    For the equivalence ``$f_n \to 0$ if and only if $1 / f_n \to \infty$'' it suffices to notice that for any compact $K \subset G$ we have
    \begin{equation*}
        \frac{1}{\sup_{z \in K} \vert f_n(z) \vert} = \inf_{z \in K} \left\vert \frac{1}{f_n(z)} \right\vert.
    \end{equation*}
    Since the $f_n$ vanish nowhere, by Hurwitz' Theorem we either have $0 \notin f(G)$ or $f \equiv 0$. In the latter case we have just shown that $1 / f_n \to \infty$ compactly.

    In the former case we have, again for any compact $K \subset G$, that $m \coloneqq \min_{z \in K} \vert f(z) \vert > 0$ and $\sup_{z \in K} \vert f_n(z) - f(z) \vert < \frac{m}{2}$ for sufficiently large $n \in \N$. Thus for all $z \in K$
    \begin{equation*}
        \frac{m}{2} > \vert f(z) - f_n(z) \vert \geq \vert f(z) \vert - \vert f_n(z) \vert \geq  m - \vert f_n(z) \vert
    \end{equation*}
    and $\vert f_n(z) \vert \geq \frac{m}{2}$. We obtain, for large $n \in \N$,
    \begin{equation*}
        \sup_{z \in K} \left\vert \frac{1}{f_n(z)} - \frac{1}{f(z)} \right\vert = \sup_{z \in K} \left\vert \frac{f(z) - f_n(z)}{f(z) f_n(z)} \right\vert \leq \sup_{z \in K} \left\vert f_n(z) - f(z) \right\vert \cdot \frac{1}{m} \cdot \frac{2}{m}
    \end{equation*}
    and therefore, after letting $n \to \infty$, that $1 / f_n \to 1 / f$ compactly in $G$.
\end{proof}

\begin{definition} \label{def:normal-family}
    Let $G \subseteq \C$ be a domain and $\F \subseteq H(G)$, then $\F$ is called:
    \begin{itemize}
        \item \emph{locally bounded}, if for every $w \in G$ there is a neighborhood $U$ of $w$ and a constant $C > 0$ such that $\vert f(z) \vert \leq C$ for all $f \in \F$ and $z \in U$.
        \item \emph{normal} in $G$, if every sequence in $\F$ has a subsequence which converges compactly in $G$ to some $f \in H(G)$. If the limit $\infty$ is also permitted it is instead called \emph{$\ast$-normal}.
    \end{itemize}
\end{definition}

The former two concepts are equivalent by the following well-known theorem:

\begin{theorem}[Montel] \label{thm:montel}
    Let $G \subseteq \C$ be a domain, then a family $\F \subseteq H(G)$ is normal if and only if it is locally bounded.
\end{theorem}

The following theorem can be interpreted as a sharpened version of Montel's Theorem and is sometimes referred to as the \emph{fundamental normality test}.

\begin{theorem} \label{lem:montel-sharpened}
    Let $G \subseteq \C$ be a domain, then any family $\F \subseteq H(G, \C \setminus \{ 0, 1 \})$ is $\ast$-normal in $G$.
\end{theorem}

\begin{proof}
    We present the proof in three steps:
    \begin{enumerate}
        \item Let $w \in G$, $c > 0$ and $\F_* \subseteq \F$ such that $\vert f(w) \vert \leq c$ for all $f \in \F_*$. We aim to show that there is an open disk at $w$ in which $\F_*$ is bounded. Select $t > 0$ such that $B_t(w) \subseteq G$. Let $f \in \F_*$, then $g(z) \coloneqq f(tz + w) \in H(\D)$. By the maximum modulus principle and Schottky's Theorem we obtain
        $$ \sup_{z \in B_{t/2}(w)} \vert f(z) \vert \leq \sup_{z \in B_{1/2}(0)} \vert g(z) \vert \leq \sup_{\vert z \vert = 1/2} \vert g(z) \vert \leq \psi(1/2, c) $$
        and $f$ is bounded on the disk $B_{t/2}(w)$. Since $f$ was arbitrary, $\F_*$ is bounded as well.

        \item Fix some $w_* \in G$ and set $\F_1 \coloneqq \{ f \in \F : \vert f(w_*) \vert \leq 1 \}$. We aim to show that $\F_1$ is locally bounded in $G$. Consider the set
        $$ U \coloneqq \{ w \in G : \F_1 \textrm{ is bounded in a neighborhood of } w \}, $$
        by (1) we have that $w_* \in U$. Note that $U$ is open in $G$, since if $\F_1$ is bounded in a disk $B_r(w)$, then for any $w' \in B_r(w)$ there is a disk $B_{r'}(w') \subseteq B_r(w)$, on which $\F_1$ is bounded as well.
        
        Assume towards a contradiction that $U \neq G$, then there exists some $w \in \partial U \cap G$ such that $\F_1$ is unbounded in every neighborhood of $w$.

        If there were some $c > 0$ such that $\vert f(w) \vert \leq c$ for all $f \in \F_1$, then by (1) there would exist an open disk centered at $w$ on which $\F_1$ would be bounded -- contradicting our assumption on $w$. Thus for every $n \in \N$, we can find some $f_n \in \F_1$ such that $\vert f_n(w) \vert \geq n$ and we obtain that $\lim_{n \to \infty} \vert f_n(w) \vert = \infty$.
        
        Set $g_n \coloneqq 1 / f_n \in \F$, then $\lim_{n \to \infty} \vert g_n(w) \vert = 0$. In particular, the family $(g_n)_{n \in \N}$ is bounded at $w$ by some constant, thus by (1) the family is bounded in some disk $B$ around $w$. By Montel's Theorem it is therefore normal in $B$, and there exists a subsequence $(g_{n_k})_{k \in \N}$ which converges compactly to a $g \in H(B)$. The functions $g_{n_k}$ have no zeros, but $g(w) = 0$; by Hurwitz's Theorem we therefore have $g \equiv 0$. Then for any $z \in B \cap U$ we have
        $$ \lim_{k \to \infty} \vert f_{n_k}(z) \vert = \lim_{k \to \infty} 1 / \vert g_{n_k}(z) \vert = \infty, $$
        contradicting the assumption that $\F_1$ is bounded in a neighborhood of such $z$. We thus have $U = G$, therefore $\F_1$ is locally bounded and by Montel's Theorem therefore normal.

        \item We can now conclude the proof. Let $(f_n)_{n \in \N}$ be a sequence in $\F$, we claim that it has some subsequence which converges compactly to some function in $H(G)$ or to $\infty$.
        
        If infinitely many $f_n$ lie in $\F_1$, then there is a subsequence $(f_{n_m})_{m \in \N}$ in $\F_1$, which by (2) has a subsequence $(f_{n_{m_k}})_{k \in \N}$ in $\F_1$ which convergences compactly in $G$ to some $f \in H(G)$. This sequence is also a subsequence of $(f_n)_{n \in \N}$, concluding the claim in this case.
        
        On the other hand, if there are only finitely many $f_n$ in $\F_1$, then infinitely many $1 / f_n$ lie in $\F_1$. As above, we thus obtain some subsequence in $\F_1$, say $(g_n)_{n \in \N}$, converging compactly in $G$ to some $g \in H(G)$. The sequence $(1 / g_n)_{n \in \N}$ is a subsequence of $(f_n)_{n \in \N}$, which -- by \Cref{lem:compact-convergence-reciprocals} -- converges compactly to $1 / g$ if $0 \notin g(G) $, and to $\infty$ otherwise. \qedhere
    \end{enumerate}
\end{proof}
\section{Picard's Great Theorem}
\label{sec:picards-great-theorem}

\begin{theorem}[Montel] \label{thm:montel}
    Let $D \subset \C$ be open. Then a family $\mathscr{F} \subset H(D)$ is normal if and only if it is locally uniformly bounded.
\end{theorem}

\begin{lemma} \label{lem:montel-sharpened}
    Let $G \subseteq \C$ be a domain and consider the family $\mathscr{F} \coloneqq H(G, \C \setminus \{ 0, 1 \})$.
    \begin{enumerate}[i.]
        \item For $z_0 \in G$ and $\omega > 0$ let $\mathscr{F}_*$ be a subfamily of $\mathscr{F}$ satisfying $\vert f(z_0) \vert \leq \omega$ for all $f \in \mathscr{F}_*$. Then there exists a neighborhood of $z_0$ in which $\mathscr{F}_*$ is bounded.
        \item For $z_0 \in G$ the family $\mathscr{F}_1 \coloneqq \{ f \in \mathscr{F} : \vert f(z_0) \vert \leq 1 \}$ is locally bounded in $G$.
        \item The family $\mathscr{F}$ is normal in $G$.
    \end{enumerate}
\end{lemma}

\begin{lemma} \label{lem:great-picard-bounded}
    Let $f \in H(\D^\times; \C \setminus \{ 0, 1 \})$. Then $f$ or $1/f$ is bounded in a punctured neighborhood of $0$.
\end{lemma}

\begin{proof}
    \todo{TODO.}
\end{proof}

\begin{theorem}[Picard's Great Theorem] \label{thm:picards-great-theorem}
    Let $G \subset \C$ be open, $w \in G$ and suppose $f \in H(G \setminus \{ w \})$ such that $f$ has an essential singularity at $w$. Then $f$ takes on all values in $\C$, with at most one exception, in any punctured neighborhood of $w$.
\end{theorem}

\begin{proof}
    Aiming for contradiction, assume that $f$ only takes on $z_0, z_1 \in \C$ finitely often in some punctured neighborhood $W$ of $w$. Then $W$ contains a punctured disk of radius $\varepsilon > 0$ around $w$, on which $f$ does not assume $z_0, z_1$. The function
    $$ g(z) \coloneqq \frac{f(\varepsilon z + w) - z_0}{z_1 - z_0} \in H(\D^\times; \C \setminus \{ 0, 1 \}) $$
    has an essential singularity at zero. By \Cref{lem:great-picard-bounded}, we have that either $g$ or $1/g$ must be bounded in a neighborhood of zero. In the former case the singularity must therefore be removable, whereas in the latter case it must be a pole, yielding a contradiction.
\end{proof}

\begin{corollary} \label{cor:transcendental-every-value-inf}
    Every entire transcendental function assumes every value in $\C$ infinitely often, with at most one exception.
\end{corollary}

\begin{proof}
    If $f(z) = \sum_{n \in \Z_{\geq 0}} a_n z^n$, then $f(z^{-1}) = \sum_{n \in \Z_{\leq 0}} a_{-n} z^{n}$, thus $f$ has an essential singularity at $\infty$ and Picard's Great Theorem (\ref{thm:picards-great-theorem}) concludes the claim.
\end{proof}

\begin{corollary}[Picard's Little Theorem] \label{thm:picards-little-theorem}
    Every nonconstant entire function omits at most one value.
\end{corollary}

\begin{proof}
    A non-constant entire function $f$ is either a non-constant polynomial or a transcendental function. In the latter case, the claim follows from \Cref{cor:transcendental-every-value-inf}.
    
    Otherwise let $w \in \C$, then $f(z) - w$ is a non-constant polynomial. By the Fundamental Theorem of Algebra it has a zero at some $z_0 \in \C$. Therefore $f(z_0) = w$ and $f$ attains all values.
\end{proof}
\chapter{Growth and Zeros of Entire Functions}
\label{ch:growth-zeros-entire-functions}

\todo{Einleitung worum es geht.}

\section{Order of Entire Functions} \label{sec:order}

To study the rate of growth of an entire function, we first introduce the following:
\begin{definition}
    Let $f \in H(\C)$, then for $r \geq 0$ we define
    \begin{equation*}
        M_f(r) \coloneqq \max_{\vert z \vert = r} \vert f(z) \vert. \qedhere
    \end{equation*}
\end{definition}

\begin{remark}    
    By the maximum modulus principle, $M_f$ is either strictly increasing (if $f$ is non-constant) or constant (otherwise), and we have
    $$ M_f(r) = \max_{\vert z \vert \leq r} \vert f(z) \vert. $$
    We shall show that $M_f$ is continuous: Let $\varepsilon > 0$ and choose $\delta_\varepsilon > 0$ such that for ${\vert z_1 - z_2 \vert < \delta_\varepsilon}$ we have ${\vert f(z_1) - f(z_2) \vert < \varepsilon}$. Now let $r_1 < r_2$ and choose $\theta$ such that $M_f(r_2) = \vert f(r_2 e^{i \theta}) \vert$. Then
    \begin{equation*}
        0 \leq M_f(r_2) - M_f(r_1) \leq \vert f(r_2 e^{i \theta}) \vert - \vert f(r_1 e^{i \theta}) \vert \leq \vert f(r_2 e^{i \theta}) - f(r_1 e^{i \theta}) \vert < \varepsilon,
    \end{equation*}
    for $r_2 - r_1 < \delta_\varepsilon$.
\end{remark}

\begin{definition} \label{def:order}
    Let $f \in H(\C)$. The \emph{order} of $f$ is defined by
    \begin{equation} \label{eq:def-order}
        \rho_f \coloneqq \limsup_{r \to \infty} \frac{\log \log M_f(r)}{\log r}.
    \end{equation}
    Constant functions, by convention, have order 0.
\end{definition}

Note that, for $f \in H(\C)$ we have $0 \leq \rho_f \leq \infty$.

\begin{proposition} \label{prop:order-exponential-polynomial}
    If $Q$ is a polynomial of degree $n \in \N$, then $\exp Q$ has order $n$.
\end{proposition}

\begin{proof}
    Write $Q(z) = \sum_{k=0}^n a_k z^k$ and let $\zeta$ be an $n$-th root of $\overline{a_n} / \vert a_n \vert$. Then
    \begin{align*}
        \vert a_n \vert - \sum_{k=0}^{n-1} \vert a_k \vert r^{k-n} \leq \Re \sum_{k=0}^{n} a_k \zeta^k r^{k-n} = \Re \frac{Q(r \zeta)}{r^n} \leq \max_{\vert z \vert = r} \Re \frac{Q(z)}{r^n} \leq \vert a_n \vert + \sum_{k=0}^{n-1} \vert a_k \vert r^{k-n}
    \end{align*}
    and taking limits as $r \to \infty$ we get $\lim_{r \to \infty} \max_{\vert z \vert = r} \Re \frac{Q(z)}{r^n} = \vert a_n \vert$. Set $f(z) \coloneqq \exp Q(z)$, then
    \begin{align*}
        \frac{\log \log M_f(r)}{\log r} &= \frac{\log \log \max_{\vert z \vert = r} \vert \exp Q(z) \vert}{\log r} = \frac{\log \log \max_{\vert z \vert = r} \exp \Re Q(z)}{\log r} = \\
        &= \frac{\log \max_{\vert z \vert = r} \Re Q(z)}{\log r} = \frac{\log \max_{\vert z \vert = r} \Re Q(z) / r^n}{\log r} + n
    \end{align*}
    and taking the limit superior as $r \to \infty$ yields $\rho_f = n$.
\end{proof}

% \begin{remark}\label{rem:order-exponential-polynomial}
%     If $Q(z) = \sum_{k=0}^n a_k z^k$ is a polynomial we claim that $f(z) \coloneqq \exp {Q(z)}$ is of order $n$. Let $\varepsilon > 0$, then since $\lim_{r \to \infty}(\sum_{k=0}^n \vert a_k \vert r^k) / r^{n + \varepsilon} = 0$ there is some $r_0$ such that for all $r > r_0$ we have $(\sum_{k=0}^n \vert a_k \vert r^k) / r^{n + \varepsilon} < 1$ and thus
%     \begin{align*}
%         M_f(r) \leq \exp \sum_{k=0}^n \vert a_n \vert r^k \leq \exp(r^{n + \varepsilon}),
%     \end{align*}
%     showing $\rho_f \leq n + \varepsilon$. On the other hand, since
% \end{remark}

If the order of an entire function is finite, we have an equivalent characterization.

\begin{proposition} \label{prop:order-infimum}
    Let $f \in H(\C)$, then $f$ is of finite order if and only if
    \begin{equation}
        \rho \coloneqq \inf \{ s > 0 : M_f(r) = \O(\exp r^s) \textrm{ as } r \to \infty \}
    \end{equation}
    is finite, and in either case we have $\rho_f = \rho$.
\end{proposition}

\begin{proof}
    Suppose $0 \leq \rho < \infty$, then for all $s > \rho$ we have $M_f(r) = \O(\exp r^s)$ as $r \to \infty$. Thus there exists a constant $K > 0$ (we may assume $K > 1$) and such that for all sufficiently large $r > 0$ we have $M_f(r) \leq K \exp r^s$. Using the fact that $\log (a+b) = \log a + \log (1 + \frac{b}{a})$ we get
    \begin{align*}
        \frac{\log \log M_f(r)}{\log r} &\leq \frac{\log (r^s + \log K)}{\log r} = s + \frac{\log (1 + (\log K) / r^s)}{\log r}. \tag{\textasteriskcentered} \label{eq:proof-order-infimum-1}
    \end{align*}
    Since
    \begin{align*}
        0 \leq \frac{\log (1 + (\log K) / r^s)}{\log r} \leq \frac{(\log K) / r^s}{\log r} = \frac{\log K}{r^s \log r} \xrightarrow{r \to \infty} 0,
    \end{align*}
    by taking the limit superior as $r \to \infty$ in \eqref{eq:proof-order-infimum-1} and then letting $s \to \rho$ we obtain $\rho_f \leq \rho$.

    Now suppose $0 \leq \rho_f < \infty$ and let $s > \rho_f$. Then, by definition of the limit superior, for all sufficiently large $r > 0$ we have $ \log \log M_f(r) \leq s \log r$ and therefore $M_f(r) \leq \exp r^s$. Thus $M_f(r) = \O(\exp r^s)$, thereby $\rho \leq s$ and letting $s \to \rho_f$ yields $\rho \leq \rho_f$.
\end{proof}

\begin{remark} \label{rem:order-zero}
    Any polynomial is of order zero. Indeed, let $Q(z) = \sum_{k=0}^n a_k z^k$ be a polynomial and $m \in \N$. Then for any $r > 1$ we have
    $$ \exp r^{1/m} = \sum_{k=0}^\infty \frac{r^{k/m}}{k!} > \frac{r^{n}}{(mn)!} $$
    and therefore
    $$ M_Q(r) = \max_{\vert z \vert = r} \vert Q(z) \vert \leq \left( \sum_{k=0}^n \vert a_k \vert \right) r^n \leq \left( (mn)! \sum_{k=0}^n \vert a_k \vert \right) \exp r^{1/m}. $$
    Since $m \in \N$ was arbitrary, \Cref{prop:order-infimum} gives $\rho_Q = 0$.

    There are, however, non-polynomial functions of order zero; one may simply construct an entire function with power series coefficients of sufficiently rapid decay. Consider
    $$ f(z) \coloneqq \sum_{k=0}^\infty \frac{z^k}{(k^2)!} $$
    and let $m \in \N$. By the above, we have $\sum_{k=0}^{m-1} \frac{r^k}{(k^2)!} \leq K \exp r^{1/m}$ for some constant $K > 0$ depending only on $m$, holding for all $r > 1$. Thus,
    $$ \sum_{k=0}^\infty \frac{r^k}{(k^2)!} \leq K \exp r^{1/m} + \sum_{k=m}^\infty \frac{r^k}{(km)!} \leq K \exp r^{1/m} + \sum_{k=0}^\infty \frac{r^{k/m}}{k!} = (K + 1) \exp r^{1/m} $$
    and we obtain $\rho_f = 0$, as above.
\end{remark}

From \Cref{prop:order-infimum}, together with some rough estimates, it follows that the order of the sum or product of two entire functions is bounded above by the larger of the two:

\begin{proposition} \label{prop:order-sum-product-estimate}
    Let $f, g \in H(\C)$ be of finite order. Then it holds that:
    \begin{enumerate}[i.]
        \item $\rho_{f + g} \leq \max \{ \rho_f, \rho_g \}$.
        \item $\rho_{fg} \leq \max \{ \rho_f, \rho_g \}$.
    \end{enumerate}
\end{proposition}

For two entire functions of different, finite order, their sum has the order of the higher of the two:

\begin{proposition} \label{prop:order-sum-equality}
    Let $f, g \in H(\C)$ be of finite order. If $\rho_f < \rho_g$, then $\rho_{f + g} = \rho_g$.
\end{proposition}

\begin{proof}
    By \Cref{prop:order-sum-product-estimate} we have 
    \begin{equation*}
        \rho_{f+g} \leq \max \{ \rho_f, \rho_g \} = \rho_g, \quad\textrm{and}\quad \rho_g = \rho_{f+g+(-f)} \leq \max \{ \rho_f, \rho_{f+g} \} = \rho_{f+g},
    \end{equation*}
    since if $\rho_f > \rho_{f+g}$ then $\rho_f > \rho_{f+g} \geq \rho_g$ would be a contradiction. Thus $\rho_{f+g} = \rho_g$, proving the assertion.
\end{proof}

It is of note that \Cref{prop:order-sum-equality} implies that the order of an entire function of finite order remains unchanged when adding a polynomial of arbitrary degree to it.

There is a, perhaps surprising, connection between the order of an entire function and the order of its derivative:

\begin{proposition} \label{prop:order-derivative}
    If $f \in H(\C)$ is of finite order, then $\rho_{f'} = \rho_f$.
\end{proposition}

\begin{proof}
    Without loss of generality we may assume $f(0) = 0$. If $f$ is a polynomial the assertion is clear. Otherwise for any $r > 0$ we have,
    \begin{align*}
        M_f(r) &= \max_{\vert z \vert = r} \vert f(z) \vert = \max_{\vert z \vert = r} \left\vert \int_{[0, z]} f'(\zeta) \diff \zeta \right\vert \leq \\
        &\leq \max_{\vert z \vert = r} \left( \vert z \vert \max_{w \in \cl{B_r(0)}} \vert f'(w) \vert \right) \leq r M_{f'}(r).
    \end{align*}
    Thus
    \begin{align*}
        \frac{\log \log M_f(r)}{\log r} &\leq \frac{\log \log r M_{f'}(r)}{\log r} = \frac{\log \log M_{f'}(r)}{\log r} + \frac{\log(1 + \log r / \log M_{f'}(r))}{\log r} \leq \\
        &\leq \frac{\log \log M_{f'}(r)}{\log r} + \frac{\log r / \log M_{f'}(r)}{\log r} = \frac{\log \log M_{f'}(r)}{\log r} + \frac{1}{\log M_{f'}(r)}
    \end{align*}
    and taking limits superior as $r \to \infty$ yields $\rho_f \leq \rho_{f'}$.

    On the other hand, Cauchy's integral formula gives
    \begin{align*}
        M_{f'}(r) &= \max_{z \in \partial B_r(0)} \vert f'(z) \vert = \max_{z \in \partial B_r(0)} \left\vert \frac{1}{2\pi i} \oint_{\partial B_1(z)} \frac{f(\zeta)}{(\zeta-z)^2} \diff \zeta \right\vert \leq \\
        &\leq \max_{\substack{z \in \partial B_r(0) \\ w \in \partial B_1(z)}} \vert f(w) \vert \leq \max_{z \in \cl{B_{r+1}(0)}} \vert f(z) \vert = M_f(r + 1).
    \end{align*}
    Therefore
    \begin{align*}
        \frac{\log \log M_{f'}(r)}{\log r} \leq \frac{\log \log M_f(r+1)}{\log (r+1)} \frac{\log (r+1)}{\log r}
    \end{align*}
    and, again, taking limits superior as $r \to \infty$ we get $\rho_{f'} \leq \rho_f$.
\end{proof}
\section{Type of Entire Functions} \label{sec:type}

For functions of finite and positive order, we can obtain a natural refinement of the concept of order:

\begin{definition}
    Let $f \in H(\C)$ be of order $0 < \rho_f < \infty$. The \emph{type} of $f$ is defined by
    \begin{equation} \label{eq:def-type}
        \tau_f \coloneqq \limsup_{r \to \infty} \frac{\log M_f(r)}{r^{\rho_f}}.
    \end{equation}

    If $\tau_f = 0$, then $f$ is said to be of \emph{minimal} type, if $0 < \tau_f < \infty$, of \emph{normal} type, and if $\tau_f = \infty$, of \emph{maximal} type.
\end{definition}

Once again, we have an equivalent characterization for functions of finite type:

\begin{proposition} \label{prop:type-infimum}
    Let $f \in H(\C)$ be of finite, positive order, then $f$ is of finite type if and only if
    \begin{equation}
        \tau \coloneqq \inf \{ t > 0 : M_f(r) = \O(\exp (t r^{\rho_f})) \textrm{ as } r \to \infty \}
    \end{equation}
    is finite, and in either case we have $\tau_f = \tau$.
\end{proposition}

\begin{proof}
    Suppose $0 \leq \tau < \infty$, then for all $t > \tau$ we have $M_f(r) = \O(\exp t r^{\rho_f})$ as $r \to \infty$. Thus there exists a constant $K > 0$ such that for all sufficiently large $r > 0$ we have $M_f(r) \leq K \exp t r^{\rho_f}$. Therefore
    \begin{align*}
        \frac{\log M_f(r)}{r^{\rho_f}} \leq \frac{\log K + t r^{\rho_f}}{r^{\rho_f}} = \frac{\log K}{r^{\rho_f}} + t
    \end{align*}
    and taking limits superior as $r \to \infty$ and letting $t \to \tau$ afterwards yields $\tau_f \leq \tau$.

    Now suppose $0 \leq \tau_f < \infty$ and let $t > \tau_f$. Then, by definition of the limit superior, for all sufficiently large $r > 0$ we have $ \log M_f(r) \leq t r^{\rho_f}$ and therefore $M_f(r) \leq \exp(t r^{\rho_f})$. Thus $M_f(r) = \O(\exp (t r^{\rho_f}))$, thereby $\tau \leq t$ and letting $t \to \tau_f$ yields $\tau \leq \tau_f$.
\end{proof}

% Equivalently, $\tau_f$ can also be defined as the infimum over all $\tau > 0$ such that $M_f(r) = \O(\exp(\tau r^{\rho_f}))$ as $r \to \infty$.

If $f \in H(\C)$ is of finite, positive order, then $f$ and $f'$ not only have the same order, they have the same type:

\begin{proposition} \label{prop:type-derivative}
    If $f \in H(\C)$ is of finite, positive order and finite type, then $\tau_{f'} = \tau_f$.
\end{proposition}

\begin{proof}
    Reusing the inequalities obtained in the proof of \Cref{prop:order-derivative} we have, for all $r > 0$,
    \begin{equation*}
        M_f(r) \leq r M_{f'}(r), \quad \textrm{and} \quad M_{f'}(r) \leq M_f(r+1).
    \end{equation*}
    Therefore
    \begin{align*}
        \frac{\log M_f(r)}{r^{\rho_f}} &\leq \frac{\log r + \log M_{f'}(r)}{r^{\rho_f}}, \quad \textrm{and} \quad \frac{\log M_{f'}(r)}{r^{\rho_f}} \leq \frac{\log M_{f}(r+1)}{(r+1)^{\rho_f}} \frac{(r+1)^{\rho_f}}{r^{\rho_f}}
    \end{align*}
    and, since by \Cref{prop:order-derivative} $\rho_f = \rho_{f'}$, taking limits superior as $r \to \infty$ yields $\tau_f \leq \tau_{f'}$ and $\tau_{f'} \leq \tau_f$.
\end{proof}

\begin{example} \label{exm:order-and-type}
    \leavevmode
    \begin{arrowlist}
        \item For $\rho, \tau \in (0, \infty)$, the function
        $$ f(z) \coloneqq \exp {\tau z^\rho} $$
        is of order $\rho$ and type $\tau$.
        \item We have already seen in \Cref{rem:order-zero} that polynomials are of order zero.
        \item The function
        $$ f(z) \coloneqq \exp \exp z $$
        is of infinite order, since for any $m \in \N$
        $$ M_f(r) \geq \exp \exp r > \exp (r^m / m!), $$
        therefore $\rho_f \geq m$, and since $m$ was arbitrary it follows that $\rho_f = \infty$.
        \qedhere
    \end{arrowlist}
\end{example}
\section{Order, Density and Exponent of Convergence of Zeros} \label{sec:zeros}

% We will mostly follow the contents of chapter 1 in Levin \cite{levin-distribution-of-zeros}.

We first recall a rather explicit connection between the moduli of the zeros of a holomorphic function and the modulus of the function itself:

\begin{theorem}[Jensen] \label{thm:jensen}
    Let $R > 0$, $f \in H(B_R(0))$ with $f(0) \neq 0$ and let $r_1, r_2, \hdots$ denote the moduli of the zeros of $f$ in $B_{R}(0)$ arranged in a non-decreasing sequence. Then, for $r_n < r < r_{n+1}$, we have
    \begin{equation}
        \frac{1}{2 \pi} \int_0^{2\pi} \log \vert f(r e^{i \theta}) \vert \diff \theta = \log \vert f(0) \vert + \log \frac{r^n}{r_1 \hdots r_n}.
    \end{equation}
\end{theorem}

\begin{definition}
    Let $f \in H(B_R(0))$ be not identically zero. Then, for $0 < r < R$, we define
    \begin{equation}
        n_f(r) \coloneqq \vert \{ z \in \cl{B_r(0)} : f(z) = 0 \} \vert,
    \end{equation}
    that is, the number of zeros of $f$ in $\cl{B_{r}(0)}$.
\end{definition}

This zero-counting function is, in some scenarios, simpler to work with than the sequence of zeros. For instance, it can be used to obtain an equivalent version of Jensen's Theorem:

\begin{corollary} \label{cor:jensen-nf}
    Let $f \in H(B_R(0))$ with $f(0) \neq 0$. Then, for $0 < r < R$, we have
    \begin{equation}
        \frac{1}{2 \pi} \int_0^{2\pi} \log \vert f(re^{i \theta}) \vert \diff \theta = \log \vert f(0) \vert + \int_0^r \frac{n_f(s)}{s} \diff s
    \end{equation}
\end{corollary}

\begin{proof}
    Let $r_1, r_2, \hdots$ denote the moduli of the zeros of $f$ in $B_{R}(0)$ arranged in a non-decreasing sequence. Then, for any $r_n < r < r_{n+1}$, we obtain
    \begin{align*}
        \log \frac{r^n}{r_1 \hdots r_n} &= \sum_{k=1}^n \log \frac{r}{r_k} = \sum_{k=1}^n \int_{r_k}^r \frac{1}{s} \diff s = \\
        &= \sum_{k=1}^n \int_0^r \1_{(r_k, \infty)}(s) \frac{1}{s} \diff s = \int_0^r \left( \sum_{k=1}^n \1_{(r_k, \infty)}(s) \right) \frac{1}{s} \diff s = \\
        &= \int_0^r \frac{n_f(s)}{s} \diff s
    \end{align*}
    and \Cref{thm:jensen} establishes the assertion.
\end{proof}

This gives an immediate connection between the zeros and the growth of the modulus of an entire function:

\begin{lemma} \label{lem:zeros-bounded-by-modulus}
    If $f \in H(\C)$, then
    $$ n_f(r) = \O( \log M_f(er) ). $$
\end{lemma}

\begin{proof}
    We may assume $\vert f(0) \vert = 1$. Let $r > 0$, then since $n_f$ is non-negative and non-decreasing we have
    \begin{equation*}
        n_f(r) \leq \int_r^{er} \frac{n_f(t)}{t} \diff t \leq \frac{1}{2 \pi} \log \vert f(er e^{i \theta}) \vert \diff \theta \leq \log M_f(er),
    \end{equation*}
    establishing the assertion.
\end{proof}

\Cref{lem:zeros-bounded-by-modulus} shows that the more zeros a function $f$ has, the faster $M_f$ must grow as $r \to \infty$. The converse is naturally false; by composing exponentials one can obtain an entire function $f$ for which $M_f$ grow arbitrarily fast, yet $f$ has no zeros.

\begin{definition} \label{def:zero-exponent}
    Let $M \subseteq \N$ and $\mathbf{z} = (z_m)_{m \in M}$ be a family of points\footnote{Note that the exponent of convergence is only affected by their moduli, not by their arguments.} in $\C^\times$. Then
    $$ \lambda_{\mathbf{z}} \coloneqq \inf \left\{ \lambda > 0 : \sum_{m \in M} \frac{1}{\vert z_m \vert^\lambda} < \infty \right\} $$
    is called the \emph{exponent of convergence} of the family $\mathbf{z}$.

    Let $f \in H(\C)$, then the \emph{exponent of convergence of the zeros of $f$} is defined as the exponent of convergence of the non-zero roots of $f$ and is denoted $\lambda_f$. If $f$ is constant, we set $\lambda_f = 0$ by convention.

    Furthermore, for any $a \in \C$, the \emph{exponent of convergence of the $a$-points of $f$} is defined as exponent of convergence of the zeros of $f(z) - a$ and will be denoted $\lambda_f^{(a)}$.
\end{definition}

If $f \in H(\C)$ assumes some value $a \in \C$ only finitely often, then clearly $\lambda_f^{(a)} = 0$.

\begin{theorem} \label{thm:inequality-order-exponent-of-convergence}
    If $f \in H(\C)$ is of finite order, then $\lambda_f \leq \rho_f$.
\end{theorem}

\begin{proof}
    We may assume $f(0) \neq 0$. If $\lambda_f = 0$ there is nothing to show. Let $\rho > \rho_f$ then
    by \Cref{lem:zeros-bounded-by-modulus} and \Cref{prop:order-infimum} we have $n_f(r) = \O(\log M_f(er))$ and $M_f(r) = \O(\exp r^\rho)$, respectively. Therefore $n_f(r) = \O(r^\rho)$, thus there exists a constant $K_1 > 0$ such that for sufficiently large $r > 0$ we have $n_f(r) \leq K_1 r^\rho$.
    
    If $(r_j)_{j \in \N}$ denote the non-zero moduli of the zeros of $f$, arranged in non-decreasing order, then for any $m \in \N$ we have $m \leq n_f(r_m) \leq K_1 r_m^{\rho}$. Let $0 < \lambda < \lambda_f$ then this implies
    \begin{align*}
        \left( \frac{1}{r_m} \right)^{\lambda} \leq K_2 \left( \frac{1}{m} \right)^{\lambda/\rho}.
    \end{align*}
    for some constant $K_2 > 0$. Therefore
    \begin{equation*}
        \sum_{m=1}^\infty \left( \frac{1}{r_m} \right)^{\lambda} \leq K_2 \sum_{m=1}^\infty \left( \frac{1}{m} \right)^{\lambda/\rho}.
    \end{equation*}
    The left-hand side diverges by \Cref{def:zero-exponent}, thus so does the right-hand side. But the latter diverges if and only if $\frac{\lambda}{\rho} \leq 1$, therefore letting $\lambda \nearrow \lambda_f$ and then $\rho \searrow \rho_f$ yields $\lambda_f \leq \rho_f$.
\end{proof}

\begin{remark}
    The function $f(z) \coloneqq \exp z$ is of order one and has no zeros -- thus we observe that we may have $\lambda_f < \rho_f$ in some cases.
\end{remark}

A more precise analysis of the density of the zeros of an entire function $f$ is contained in the growth of the zero-counting function $n_f$. We therefore introduce the following definitions:

\begin{definition} \label{def:order-density-zero-counting}
    Let $f \in H(\C)$. As the \emph{order} of the zero-counting function $n_f$ we define
    \begin{equation}
        \nu_f \coloneqq \limsup_{r \to \infty} \frac{\log n_f(r)}{\log r}.
    \end{equation}
    If $\nu_f$ is finite, we additionally define the \emph{upper density} of the zeros of $f$ as
    \begin{equation}
        \Delta_f \coloneqq \limsup_{r \to \infty} \frac{n_f(r)}{r^{\nu_f}}.
    \end{equation}
    If the limit exists, then $\Delta_f$ is simply called the \emph{density}.
\end{definition}

The following lemma provides great utility in further analyzing the function $n_f$.

\begin{lemma} \label{lem:zeros-measure}
    Let $f \in H(\C)$ and denote by $(r_m)_{m \in M}$ the moduli of the zeros of $f$ arranged in non-decreasing order. Then $f$ induces a unique -Stieltjes measure $\omega_f$ on $(0, \infty)$, such that, for any non-negative, measureable function $\varphi$ on $(0, \infty)$, it holds that
    \begin{equation*}
        \sum_{m \in M} \varphi(r_m) = \int_{(0, \infty)} \varphi \diff \omega_f.
    \end{equation*}
    If in addition $\varphi$ is monotone and continuously differentiable on $(0, \infty)$, then for all $r > 0$ it holds that
    \begin{equation*}
        \int_{(0, r)} \varphi \diff \omega_f = n_f(r) \varphi(r) - \int_0^r n_f(t) \varphi'(t) \diff t.
    \end{equation*}
\end{lemma}

\begin{proof}
    From the representation
    $$ n_f(r) = \sum_{m \in M} \1_{[r_m, \infty)}(r) $$
    we observe that $n_f$ is right-continuous, non-negative and non-decreasing. Therefore $n_f$ induces a unique Lebesgue--Stieltjes measure $\omega_f$. Since an indicator function of the form $\1_{[a, \infty)}$, for $a \in \R$, corresponds to a Dirac measure $\delta_{\{ a \}}$, and the correspondence between Lebesgue-Stieltjes measures and their distribution functions is linear in nature, we have
    $$ \omega = \sum_{m \in M} \delta_{\{r_m\}}. $$
    From this it follows that
    \begin{equation*}
        \sum_{m \in M} \varphi(r_m) = \sum_{m \in M} \int_{(0, \infty)} \varphi \diff \delta_{\{r_m\}} = \int_{(0, \infty)} \varphi \diff \left( \sum_{m \in M} \delta_{\{r_m\}} \right) = \int_{(0, \infty)} \varphi \diff \omega_f.
    \end{equation*}
    Assume $\varphi$ is monotone and continuously differentiable on $(0, \infty)$. Then $\varphi$ also induces a Lebesgue--Stieltjes measure $\eta_\varphi$, with density $\varphi'$ with respect to the Lebesgue measure. Thus, for any $0 < \varepsilon < r < \infty$, integration by parts gives
    \begin{equation*}
        \int_{(\varepsilon, r)} \varphi \diff \omega_f + \int_{(\varepsilon, r)} n_f \diff \eta_\varphi = n_f(r) \varphi(r) - n_f(\varepsilon) \varphi(\varepsilon).
    \end{equation*}
    Since $\kappa \coloneqq \min_{m \in M} r_m > 0$ we have $n_f(s) = 0$ for $0 < s < \kappa$. Letting $\varepsilon \to 0$ and using the density of $\eta_\varphi$ therefore yields
    \begin{equation*}
        \int_{(0, r)} \varphi \diff \omega_f + \int_0^r n_f(t) \varphi'(t) \diff t = n_f(r) \varphi(r)
    \end{equation*}
    and rearranging terms provides the desired result.
\end{proof}

\begin{proposition} \label{prop:zeros-order-equals-zeros-exponent}
    Let $f \in H(\C)$, then $\lambda_f$ is finite if and only if $\nu_f$ is finite and in either case we have $\lambda_f = \nu_f$.
\end{proposition}

\begin{proof}
    If $f$ only has finitely many zeros the assertion is clear. Let $\lambda > \lambda_f$ and $r > 0$. Invoking \Cref{lem:zeros-measure} with $\varphi(t) \coloneqq t^{-\lambda}$ we get
    \begin{equation*}
        \int_{(0, r)} \frac{1}{t^\lambda} \diff \omega_f(t) = \frac{n_f(r)}{r^\lambda} + \lambda \int_0^r \frac{n_f(t)}{t^{\lambda + 1}} \diff t.
    \end{equation*}
    Since $\lambda > \lambda_f$, the same lemma implies that the left integral converges as $r \to \infty$. In particular, the right-hand side must be bounded. Since the right integrand is non-negative, the right-most term is increasing in $r$ and must therefore converge as $r \to \infty$. Therefore the right-most term in
    $$ 0 \leq \frac{n_f(r)}{r^\lambda} = \lambda n_f(r) \int_r^\infty \frac{1}{t^{\lambda + 1}} \diff t \leq \lambda \int_r^\infty \frac{n_f(t)}{t^{\lambda + 1}} \diff t $$
    converges to $0$ as $r \to \infty$, from which we get
    \begin{equation} \label{eq:zeros-order-equals-zeros-exponent:density-zero}
        \lim_{r \to \infty} \frac{n_f(r)}{r^\lambda} = 0.
    \end{equation}
    Taking limits superior as $r \to \infty$ in
    $$ \frac{\log n_f(r)}{\log r} = \left( \log \frac{n_f(r)}{r^\lambda} + \log r^\lambda\right)/{\log r} = \log \frac{n_f(r)}{r^\lambda} / \log r + \lambda $$
    and then letting $\lambda \searrow \lambda_f$ yields $\nu_f \leq \lambda_f$.

    Conversely suppose $\nu_f < \infty$ and let $\varepsilon > 0$. By the definition of $\nu_f$ and the nature of the limit superior there is an $r_0 > 0$ such that
    $$ n_f(r) \leq r^{\nu_f + \varepsilon} $$
    for all $r \geq r_0$. Set $\lambda \coloneqq \nu_f + 2 \varepsilon$, then
    \begin{equation*}
        \int_{r_0}^r \frac{n_f(t)}{t^{\lambda + 1}} \diff t \leq
        \int_{r_0}^r \frac{t^{\nu_f + \varepsilon}}{t^{\nu_f + 2\varepsilon + 1}} \diff t =
        \int_{r_0}^r t^{-\varepsilon - 1} \diff t = \frac{1}{\varepsilon r_0^{\varepsilon}} - \frac{1}{\varepsilon r^{\varepsilon}}
    \end{equation*}
    and since the right-hand side remains finite as $r \to \infty$ therefore the leftmost integral converges as $r \to \infty$. Clearly \eqref{eq:zeros-order-equals-zeros-exponent:density-zero} holds for our $\lambda$, therefore \Cref{lem:zeros-measure} implies the convergence of
    \begin{equation*}
        \int_{(0, \infty)} \frac{1}{t^\lambda} \diff \omega_f(t) = \sum_{m \in M} \frac{1}{r_m^\lambda},
    \end{equation*}
    yielding $\lambda_f \leq \lambda$, and letting $\varepsilon \to 0$ results in $\lambda_f \leq \nu_f$.
\end{proof}

\iffalse
\begin{example} \label{exm:exponent-of-convergence}
    Consider $f(z) \coloneqq \sin(z) \in H(\C)$, we want to calculate $\lambda_f$ and $\rho_f$. First, let $\lambda > 0$ and recall that $f$ has zeros at $(n \pi)_{n \in \Z}$. Since
    $$ \sum_{n \in \Z \setminus \{0\}} \frac{1}{\vert n \pi \vert^\lambda} = \frac{2}{\pi^\lambda} \sum_{n=1}^\infty \frac{1}{n^\lambda} $$
    is finite if and only if $\lambda > 1$, we obtain $\lambda_f = 1$. Furthermore, we have $\vert \sin(z) \vert \leq e^{\vert z \vert}$
    and therefore $M_f(r) \leq e^r$, whereby $\rho_f \leq 1$. Finally, \Cref{thm:inequality-order-exponent-of-convergence} concludes $\rho_f = 1$.
\end{example}
\fi
\section{Canonical Products} \label{sec:canonical-products}

By the Fundamental Theorem of Algebra every complex polynomial can be written as a product of linear factors involving only its roots and some scaling factor. Conversely, for a finite sequence of points $z_1, \hdots, z_n$ the polynomial $p(z) \coloneqq \prod_{k=1}^n (z - z_k)$ vanishes only at said points.

The question arises if one can construct an entire function that vanishes precisely at the points of some infinite sequence $(z_k)_{k \in \N}$ and nowhere else. Clearly the product $\prod_{k=1}^\infty (z - z_k)$ need not be convergent in the entire complex plane. This issue is solved by introducing additional exponential factors, which improve convergence without introducing new zeros. We first study sufficient conditions for convergence of infinite products of numbers and holomorphic functions.

The early parts of the section will closely follow Stein and Shakarchi \cite{stein-shakarchi-princeton}, while the later estimates and results will follow Levin \cite{levin-distribution-of-zeros}.

\begin{lemma} \label{lem:infinite-product-criteria}
    Suppose $G \subseteq \C$ is open and let $(f_k)_{k \in \N}$ be a sequence of non-vanishing functions in $H(G)$. If $\sum_{k=1}^\infty (1 - f_k(z))$ converges absolutely compactly in $G$, then $\prod_{k=1}^\infty f_k(z)$ converges compactly in $G$ to a non-vanishing $f \in H(G)$.
\end{lemma}
    
\begin{proof}
    Let $K \subset G$ be compact, then the absolute, compact convergence of the sum implies that there is a $k_0$ such that $\vert 1 - f_k(z) \vert \leq 1/2$ for $z \in K, k > k_0$. For $\vert w \vert \leq 1/2$ we have
    \begin{equation*}
        \vert \log(1 - w) \vert \leq \sum_{k=1}^\infty \left\vert \frac{1}{k} w^k \right\vert \leq \vert w \vert \sum_{k=0}^\infty \vert w \vert^k = \frac{\vert w \vert}{1 - \vert w \vert} \leq 2 \vert w \vert.
    \end{equation*}
    Therefore
    \begin{equation*}
        \sum_{k=1}^\infty \vert \log f_k(z) \vert = \sum_{k=1}^\infty \vert \log (1 - (1 - f_k(z))) \vert \leq \sum_{k=1}^{k_0 - 1} \vert \log f_k(z) \vert + 2 \sum_{k=k_0}^\infty \vert 1 - f_k(z) \vert
    \end{equation*}
    and since the rightmost sum converges compactly it follows that $\sum_{k=1}^\infty \log f_k(z)$ converges absolutely compactly in $G$. Since the exponential function is uniformly continuous, the sequence
    $$ \exp \left( \sum_{k=k_0}^n \log f_k(z) \right) = \prod_{k=k_0}^n f_k(z), \quad n \in \N $$
    converges absolutely compactly to some $f \in H(G)$. Again by continuity of the exponential function we have
    $$ f(z) = \prod_{k=1}^\infty f_k(z) = \exp \left( \sum_{k=1}^\infty \log f_k(z) \right) \neq 0 $$
    and thus $f$ is non-vanishing.
\end{proof}

\begin{definition} \label{def:canonical-factors}
    The \emph{canonical factors} are defined at $z \in \C$ by
    $$ W(z; 0) \coloneqq 1-z, \quad \textrm{and} \quad W(z; n) \coloneqq (1-z) \exp \left( \sum_{k=1}^n \frac{z^k}{k} \right), \quad n \in \N. $$
    The number $n$ is called the \emph{degree} of the canonical factor.
\end{definition}

\begin{lemma} \label{lem:estimate-canonical-factors}
    Let $n \in \N_0$ and $\vert z \vert \leq 1/2$, then $ \vert 1 - W(z; n) \vert \leq 2e \vert z \vert^{n+1} $.
\end{lemma}

\begin{proof}
    For $n = 0$ there is nothing to show. Note that for such given $z$ we can expand $\log (1 - z) = - \sum_{k=1}^\infty \frac{z^k}{k}$, therefore
    \begin{equation*}
        W(z; n) = \exp \left( \log(1-z) + \sum_{k=1}^n \frac{z^k}{k} \right) = \exp \zeta,
    \end{equation*}
    where $\zeta \coloneqq - \sum_{k=n+1}^\infty z^k / k$. Furthermore, we have that
    \begin{equation*}
        \vert \zeta \vert \leq \vert z \vert^{n+1} \sum_{k=n+1}^\infty \vert z \vert^{k-n-1} / k \leq \vert z \vert^{n+1} \sum_{j=0}^\infty 2^{-j} \leq 2 \vert z \vert^{n+1}
    \end{equation*}
    and in particular $\vert \zeta \vert \leq 1$. We conclude
    \begin{equation*}
        \vert 1 - W(z; n) \vert = \vert \exp \zeta - 1 \vert \leq \sum_{k=1}^\infty \frac{\vert \zeta \vert^k}{k!} \leq \vert \zeta \vert \sum_{k=0}^\infty \frac{1}{k!} \leq 2 e \vert z \vert^{n+1}.
    \end{equation*}
\end{proof}

The next theorem -- which is also known as the \emph{Weierstraß Product Theorem} -- now constructively asserts the existence of entire function with prescribed zeros.

\begin{theorem} \label{thm:function-with-prescribed-zeros}
    Let $M \subseteq \N$ and $\mathbf{z} = (z_m)_{m \in M}$ be a family in $\C^\times$ without accumulation points. Then there is a family $(p_m)_{m \in M}$ in $\N_0$ such that
    \begin{equation}
        \Pi(z) \coloneqq \prod_{m \in M} W\left(\frac{z}{z_m}; p_m\right)
    \end{equation}
    is an entire function that vanishes at all $z_m, m \in M$ and nowhere else. If
    \begin{equation}
        \mu \coloneqq \min \left\{ p \in \N_0 : \sum_{m \in M} \frac{1}{\vert z_m \vert^{p+1}} < \infty \right\}
    \end{equation}
    is finite, then one can take $p_m = \mu, m \in M$.
\end{theorem}

\begin{proof}
    If $M$ is finite there is nothing to show. We may thus assume that $M$ is infinite and without loss of generality $M = \N$. Thus $\mathbf{z}$ is a sequence, which we can assume to be arranged in non-decreasing order according to the moduli its elements.

    We claim that there exists a sequence $(p_n)_{n \in \N}$ in $\N_0$ such that
    \begin{equation}
        \sum_{n=1}^\infty \left\vert \frac{z}{z_n} \right\vert^{p_n + 1}
    \end{equation}
    converges compactly in $\C$. Let $K \subset B_R(0) \subset \C$ be compact. Note that since $\mathbf{z}$ has no accumulation points we necessarily have $\lim_{n \to \infty} \vert z_n \vert = \infty$, thus there is some $n_0$ such that for $n > n_0$ we have $\vert z_n \vert \geq R + 1$. Therefore, for any $z \in K$,
    \begin{equation*}
        \sum_{n=1}^\infty \left\vert \frac{z}{z_n} \right\vert^{p_n + 1} \leq \sum_{n=1}^{n_0 - 1}  \left\vert \frac{z}{z_n} \right\vert^{p_n + 1} + \sum_{n=n_0}^\infty \left( \frac{R}{R+1} \right)^{p_n + 1},
    \end{equation*}
    and the rightmost series is finite if, for example, $p_n = n, n \in \N$. Therefore the series is uniformly bounded and thus compactly convergent, concluding our claim. Note that if $\mu$ is finite then we can take $p_n = \mu, n \in \N$, since for $z \in K$
    \begin{equation*}
        \sum_{n=1}^\infty \left\vert \frac{z}{z_n} \right\vert^{\mu + 1} \leq \vert z \vert^\mu \sum_{n=1}^\infty \frac{1}{\vert z_n \vert^{\mu+1}} \leq R^\mu \sum_{n=1}^\infty \frac{1}{\vert z_n \vert^{\mu+1}},
    \end{equation*}
    where the rightmost sum is a finite constant by the definition of $\mu$.

    It remains to show that is $\Pi$ is an entire function that satisfies the desired properties.Let $R > 0$, then it suffices to show that $\Pi$ has is holomorphic in $B_R(0)$ and vanishes only at points of $\mathbf{z}$ that are also in $B_R(0)$. We write
    $$ \Pi(z) = \left( \prod_{\substack{n=1 \\ \vert z_n \vert < 2R }}^\infty W\left(\frac{z}{z_n}; p_n\right) \right) \left( \prod_{\substack{n=1 \\ \vert z_n \vert \geq 2R }}^\infty W\left(\frac{z}{z_n}; p_n\right) \right) \eqqcolon \Pi_1(z) \Pi_2(z). $$
    Clearly $\Pi_1$ is a finite product which vanishes at points of $\mathbf{z}$ in $B_R(0)$ and nowhere else. Note that for $\vert z_m \vert \geq 2R$ we have $\vert z / z_m \vert \leq R / \vert z_m \vert \leq 1 / 2$, hence by \Cref{lem:estimate-canonical-factors} we have
    $$ \sum_{\substack{n=1 \\ \vert z_n \vert \geq 2R }}^\infty \left\vert 1 - W\left(\frac{z}{z_n}; p_n\right) \right\vert \leq 2e \sum_{\substack{n=1 \\ \vert z_n \vert \geq 2R }}^\infty \left\vert \frac{z}{z_n} \right\vert^{p_n + 1} $$
    and the compact convergence of the right sum, which stems from our choice of the $p_n, n \in \N$, implies compact convergence of the left sum. By \Cref{lem:infinite-product-criteria} we thereby have that $\Pi_2 \in H(B_R(0))$ and vanishes nowhere. Therefore $\Pi$ has the desired properties and, since $R$ was arbitrary, this concludes the proof.
\end{proof}

\begin{definition} \label{def:canonical-product}
    With the notation of \Cref{thm:function-with-prescribed-zeros}, $\Pi$ is called a \emph{Weierstraß product}\footnote{Note that this product is not uniquely determined, since the sequence $(p_m)_{m \in M}$ is not unique.} formed from the family $\mathbf{z}$. If $\mu < \infty$, then
    \begin{equation}
        \Pi(z) \coloneqq \prod_{m \in M} W\left(\frac{z}{z_m}; \mu \right)
    \end{equation}
    is called a \emph{(Weierstraß) canonical product} and $\mu$ is called the \emph{genus} of the canonical product, which will also be denoted as $\mu_\Pi$. Otherwise the genus of $\Pi$ is said to be infinite.
\end{definition}

\begin{remark}
    By comparing \Cref{def:zero-exponent} and \Cref{def:canonical-product} we immediately observe that if $\lambda_\Pi < \infty$, then $\mu_\Pi \leq \lambda_\Pi \leq \mu_\Pi + 1$. If therefore $\lambda_\Pi$ is an integer, and $\Pi$ has infinitely many zeros, the series
    $$ \sum_{m \in M} \frac{1}{\vert z_m \vert^{\lambda}}, $$
    will converge if $\lambda_\Pi = \mu_\Pi + 1$ and diverge if $\lambda_\Pi = \mu_\Pi$.
\end{remark}

Just as any polynomial decomposes into linear factors involving their roots by the Fundamental Theorem of Algebra, Weierstraß products allow us to obtain a similar factorization result for arbitrary entire functions, known as the \emph{Weierstraß Factorization Theorem}.

\begin{theorem}[Weierstraß] \label{thm:Weierstraß}
    Let $f \in H(\C)$ be not identically zero and let $m \in \N_0$ be the order of the zero of $f$ at the origin. Then there exists $g \in H(\C)$ such that
    \begin{equation}
        f(z) = z^m e^{g(z)} \Pi(z),
    \end{equation}
    where $\Pi$ denotes a Weierstraß product formed from the zeros of $f$.
\end{theorem}

\begin{proof}
    Since $f(z)$ and $z^m \Pi(z)$ have the same zeros with identical multiplicities the function
    $$ \varphi(z) \coloneqq \frac{f(z)}{z^m \Pi(z)} $$
    is entire and vanishes nowhere. Therefore we can write
    $$ \varphi(z) = e^{g(z)} $$
    for some $g \in H(\C)$ and rearranging terms yields the desired representation.
\end{proof}

\begin{definition} \label{def:genus}
    In the context of Weierstraß' Theorem, if $\Pi$ is of finite genus and $g$ is a polynomial, then the \emph{genus of the entire function $f$} is defined as
    \begin{equation}
        \mu_f \coloneqq \max \{ \mu_\Pi, \deg g \}.
    \end{equation}

    If $g$ is not a polynomial or if $\Pi$ is of infinite genus, then $f$ is said to be of infinite genus.
\end{definition}

We wish to further estimate canonical products. To do this we first focus on the canonical factors.

\begin{lemma} \label{lem:elementary-factor-estimate}
    For $p \in \N$ and all $z \in \C$ we have
    \begin{equation}
        \log \vert W(z; p) \vert < A_p \frac{\vert z \vert^{p+1}}{1 + \vert z \vert}, \quad \textrm{where} \quad A_p \coloneqq 3e (2 + \log p).
    \end{equation}
    For $p = 0$ we have
    \begin{equation}
        \log \vert W(z; 0) \vert \leq \log (1 + \vert z \vert).
    \end{equation}
\end{lemma}

\begin{proof}
    The second assertion is clear since
    $$ \log \vert W(z; 0) \vert = \log \vert 1 - z \vert \leq \log (1 + \vert z \vert). $$
    If $\vert z \vert \leq \frac{p}{p+1}$ then in particular we have $\vert z \vert < 1$. Since furthermore $1 + \vert z \vert < 2 < A_p$ we get
    \begin{align*}
        \log \vert W(z; p) \vert &= \log \left\vert (1-u) \exp \left( \sum_{k=1}^p \frac{z^k}{k} \right) \right\vert = \Re \left( \log(1-u) + \sum_{k=1}^p \frac{z^k}{k} \right) = \\
        &= -\Re \sum_{k=p+1}^\infty \frac{z^k}{k} \leq \sum_{k=p+1}^\infty \frac{\vert z \vert^k}{k} < \frac{\vert z \vert^{p+1}}{(p+1)(1 - \vert z \vert)} \leq \vert z \vert^{p+1} \leq \\ &\leq A_p \frac{\vert z \vert^{p+1}}{1 + \vert z \vert}.
    \end{align*}
    If otherwise $\vert z \vert > \frac{p}{p+1}$ then we have $\frac{1}{\vert z \vert} < 1 + \frac{1}{p}$ and $2 \vert z \vert > 1$. From there we conclude
    \begin{align*}
        \log \vert W(z; p) \vert &= \log \left\vert (1-z) \exp \left( \sum_{k=1}^p \frac{z^k}{k} \right) \right\vert \leq \log(1 + \vert z \vert) + \sum_{k=1}^p \frac{\vert z \vert^k}{k} \leq \\
        &\leq 2 \vert z \vert + \sum_{k=2}^p \frac{\vert z \vert^k}{k} = \vert z \vert^p \left( 2 \left\vert \frac{1}{z} \right\vert^{p-1} + \sum_{k=2}^p \frac{1}{k} \left\vert \frac{1}{z} \right\vert^{p-k} \right) < \\
        &< \vert z \vert^p \left( 2 \left(1 + \frac{1}{p} \right)^{p-1} + \sum_{k=2}^p \frac{1}{k} \left( 1 + \frac{1}{p} \right)^{p-k} \right) < \\
        &< \vert z \vert^p \left( 1 + \frac{1}{p} \right)^p \left( 2 + \sum_{k=1}^p \frac{1}{k} \right) < \vert z \vert^p e \left( 2 + \int_1^p \frac{1}{t} \diff t \right) = \\
        &= \vert z \vert^p e \left( 2 + \log p \right) = \frac{1 + \vert z \vert}{1 + \vert z \vert} \vert z \vert^p e \left( 2 + \log p \right) < \\
        &< \frac{3 \vert z \vert}{1 + \vert z \vert} \vert z \vert^p e \left( 2 + \log p \right) = A_p \frac{\vert z \vert^{p+1}}{1 + \vert z \vert}. \qedhere
    \end{align*}
\end{proof}

\begin{lemma} \label{lem:canonical-product-estimate}
    Let $M \subseteq \N$ and $(z_m)_{m \in M}$ be a family in $\C^\times$ without accumulation points. If for given $p \in \N_0$ it holds that
    $$ \sum_{m \in M} \frac{1}{\vert z_m \vert^{p+1}} < \infty, $$
    then the entire function
    $$ \Pi(z) \coloneqq \prod_{m \in M} W \left( \frac{z}{z_m}; p \right) $$
    satisfies
    \begin{equation}
        \log \vert \Pi(z) \vert < k_p r^p \left( \int_0^r \frac{n_\Pi(t)}{t^{p+1}} \diff t + r \int_r^\infty \frac{n_\Pi(t)}{t^{p+2}} \diff t \right),
    \end{equation}
    for all $z \in \C^\times$ and $r = \vert z \vert$, where
    $$ k_0 \coloneqq 1, \quad \textrm{and} \quad k_p \coloneqq 3e(p+1)(2 + \log p), \quad p \in \N. $$
\end{lemma}

\begin{proof}
    Assuming $p > 0$, combining \Cref{lem:elementary-factor-estimate,lem:zeros-measure} gives
    \begin{align*}
        \log \vert \Pi(z) \vert &= \prod_{m \in M} W \left( \frac{z}{z_m}; p \right) = \sum_{m \in M} \log \left\vert W \left( \frac{z}{z_m}; p \right) \right\vert < \\
        &< A_p \sum_{m \in M} \frac{\left\vert \frac{z}{z_m} \right\vert^{p+1}}{1 + \left\vert \frac{z}{z_m} \right\vert} = A_p r^{p+1} \sum_{m \in M} \frac{1}{\vert z_m \vert^p (\vert z_m \vert + r)} = \\
        &= A_p r^{p+1} \int_{(0, \infty)} \frac{1}{t^p (t+r)} \diff \omega_\Pi(t),
    \end{align*}
    \Cref{lem:zeros-measure} also yields, for $s > 0$,
    \begin{equation*}
        \int_{(0, s)} \frac{1}{t^p (t+r)} \diff \omega_\Pi(t) = \frac{n_\Pi(s)}{s^p (s + r)} + \int_{(0, s)} n_\Pi(t) \frac{(p+1)t + pr}{t^{p+1} (t+r)^2} \diff t.
    \end{equation*}
    As seen in the proof of \Cref{prop:zeros-order-equals-zeros-exponent}, the left term on the right-hand side tends to $0$ as $s \to \infty$. Therefore we obtain
    \begin{align*}
        \log \vert \Pi(z) \vert &< A_p r^{p+1} \int_{(0, \infty)} \frac{1}{t^p (t+r)} \diff \omega_\Pi(t) = A_p r^{p+1} \int_0^\infty n_\Pi(t) \frac{(p+1)t + pr}{t^{p+1} (t+r)^2} \diff t < \\
        &< A_p (p+1) r^{p+1} \int_0^\infty \frac{n_\Pi(t)}{t^{p+1}(t+r)} \diff t < \\
        &< A_p (p+1) r^{p+1} \left( \frac{1}{r} \int_0^r \frac{n_\Pi(t)}{t^{p+1}} \diff t + \int_r^\infty \frac{n_\Pi(t)}{t^{p+2}} \diff t \right) = \\
        &= k_p r^{p} \left( \int_0^r \frac{n_\Pi(t)}{t^{p+1}} \diff t + r \int_r^\infty \frac{n_\Pi(t)}{t^{p+2}} \diff t \right).
    \end{align*}
    The case $p=0$ is verified in the same fashion:
    \begin{align*}
        \log \vert \Pi(z) \vert &= \sum_{m \in M} \log \left\vert W \left( \frac{z}{z_m}; 0 \right) \right\vert \leq \sum_{m \in M} \log \left( 1 + \left\vert \frac{z}{z_m} \right\vert \right) \leq \\
        &\leq \sum_{m \in M} \left\vert \frac{z}{z_m} \right\vert = r \sum_{m \in M} \frac{1}{\vert z_m \vert} = r \int_{(0, \infty)} \frac{1}{t} \diff \omega(t) = -r \int_0^\infty -\frac{n_\Pi(t)}{t^2} \diff t < \\
        &< r \left( \int_0^r \frac{n_\Pi(t)}{tr} \diff t + \int_r^\infty \frac{n_\Pi(t)}{t^2} \diff t \right) = \int_0^r \frac{n_\Pi(t)}{t} \diff t + r \int_r^\infty \frac{n_\Pi(t)}{t^2} \diff t. \qedhere
    \end{align*}
\end{proof}

\begin{theorem} \label{thm:exponent-of-convergence-Weierstraß-product}
    Let $\Pi \in H(\C)$ be a canonical product of finite order, then $\lambda_\Pi = \rho_\Pi$.
\end{theorem}

\begin{proof}
    By \Cref{thm:inequality-order-exponent-of-convergence} we already have $\lambda_\Pi \leq \rho_\Pi$, it remains to show that $\rho_\Pi \leq \lambda_\Pi$. We first recall that the convergence exponent always satisfies $\mu \leq \lambda_\Pi \leq \mu + 1$, where $\mu \coloneqq \mu_\Pi$.
    
    Suppose $\lambda_\Pi < \mu + 1$ and let $\lambda_\Pi < \lambda < \mu + 1$. By \Cref{prop:zeros-order-equals-zeros-exponent} we have
    $$ \lambda_\Pi = \nu_\Pi = \limsup_{r \to \infty} \frac{\log n_\Pi(r)}{\log r} < \lambda. $$
    Therefore there exists some $r_0 > 0$ such that for all $r \geq r_0$ we have $n_\Pi(r) \leq r^\lambda$. Choose $\varepsilon > 0$ small enough such that $n_\Pi(\varepsilon) = 0$, then
    \begin{equation*}
        n_\Pi(r) = \frac{n_\Pi(r)}{r^\lambda} r^\lambda \leq \frac{n_\Pi(r_0)}{\varepsilon^\lambda} r^\lambda
    \end{equation*}
    and choosing $C_\lambda > \max \{ 1, n_\Pi(r_0) / \varepsilon^\lambda \}$ yields $n_\Pi(r) \leq C_\lambda r^\lambda$ for all $r > 0$. Applying this inequality to \Cref{lem:canonical-product-estimate} we get, for $r = \vert z \vert$,
    \begin{align*}
        \log \vert \Pi(z) \vert &< k_\mu r^\mu \left( \int_0^r \frac{n_\Pi(t)}{t^{\mu + 1}} \diff t + r \int_r^\infty \frac{n_\Pi(t)}{t^{\mu + 2}} \diff t \right) \leq \\
        &\leq k_\mu r^\mu C_\lambda \left( \int_0^r t^{\lambda - \mu - 1} \diff t + r \int_r^\infty t^{\lambda - \mu - 2} \diff t \right) = \\
        &= k_\mu r^\mu C_\lambda \left( \frac{r^{\lambda - \mu}}{\lambda - \mu} - \frac{r^{\lambda - \mu}}{\lambda - \mu - 1} \right) = k_\mu C_\lambda \left( \frac{1}{\lambda - \mu} - \frac{1}{\lambda - \mu - 1} \right) r^\lambda.
    \end{align*}
    From this it follows $M_\Pi(r) = \O(\exp r^\lambda)$ and thus that $\rho_\Pi \leq \lambda$. Letting $\lambda \searrow \lambda_\Pi$ yields $\rho_\Pi \leq \lambda_\Pi$.

    If we otherwise assume $\lambda_\Pi = \mu + 1$, then by \Cref{lem:zeros-measure}
    $$ \int_{(0, r)} \frac{1}{t^{\mu + 1}} \diff \omega_\Pi(t) = \frac{n_\Pi(r)}{r^{\mu + 1}} + (\mu + 1) \int_0^r \frac{n_\Pi(t)}{t^{\mu + 2}} \diff t, $$
    wherein the leftmost integral converges as $r \to \infty$. Imitating the proof of \Cref{prop:zeros-order-equals-zeros-exponent} we see that the left term on the right-hand side converges to $0$ as $r \to \infty$, thus the rightmost integral converges as well. Let $\varepsilon > 0$. These two convergences imply that there is an $r_0 > 0$ such that for all $r \geq r_0$ we have
    $$ \int_{r_0}^\infty \frac{n_\Pi(t)}{t^{\mu + 2}} \diff t < \varepsilon, \quad \textrm{and} \quad \frac{n_\Pi(r)}{r^{\mu + 1}} < \varepsilon. $$
    Inserting these results into \Cref{lem:canonical-product-estimate} yields, for $r = \vert z \vert \geq r_0$,
    \begin{align*}
        \log \vert \Pi(z) \vert &< k_\mu r^\mu \left( \int_0^r \frac{n_\Pi(t)}{t^{\mu + 1}} \diff t + r \int_r^\infty \frac{n_\Pi(t)}{t^{\mu + 2}} \diff t \right) = \\
        &= k_\mu r^\mu \left( \int_0^{r_1} \frac{n_\Pi(t)}{t^{\mu + 1}} \diff t + \int_{r_1}^r \frac{n_\Pi(t)}{t^{\mu + 1}} \diff t + r \int_r^\infty \frac{n_\Pi(t)}{t^{\mu + 2}} \diff t \right) \leq \\
        &\leq k_\mu r^\mu \left(C + (r - r_1) \varepsilon + r \varepsilon \right) \leq k_\mu (C r^\mu + 2 \varepsilon r^{\mu + 1}).
    \end{align*}
    The right-most term asymptotically dominates, therefore $M_\Pi(r) = \O(\exp(\varepsilon r^{\mu + 1}))$ and thus $\Pi$ is at most of order $\mu + 1 = \lambda_\Pi$ and, since $\varepsilon$ was arbitrary, minimal type.
\end{proof}

Just as the order of a canonical product is determined by the order of growth of its zeros, its type is (to some extent) determined by the upper density of its zeros. A peculiarity arises when considering canonical products of integral order, in which case the type may also depend on the argument of the zeros \cite{levin-distribution-of-zeros}. However, for non-integral order we can obtain the following result:

\begin{theorem} \label{thm:canonical-product-type-density}
    Let $\Pi \in H(\C)$ be a canonical product of finite order. If $\lambda_\Pi$ is not an integer, then $\Pi$ is of maximal, minimal or normal type according to whether $\Delta_\Pi$ is equal to infinity, zero, or a positive, real number.
\end{theorem}

\begin{proof}
    by \Cref{lem:zeros-bounded-by-modulus} there is a constant $K > 0$ such that for all sufficiently large $r$ we have $n_\Pi(r) \leq K \log M_\Pi(er)$ and by \Cref{prop:zeros-order-equals-zeros-exponent} and \Cref{thm:exponent-of-convergence-Weierstraß-product} we have $\nu_\Pi = \lambda_\Pi = \rho_\Pi$. Therefore we have, for large $r$,
    \begin{equation} \label{eq:density-smaller-constant-type}
        \frac{n_\Pi(r)}{r^{\nu_\Pi}} = \frac{n_\Pi(r)}{r^{\rho_\Pi}} \leq K e^{\rho_\Pi} \frac{\log M_\Pi(er)}{(e r)^{\rho_\Pi}}.
    \end{equation}
    We now consider two cases. If $\Delta_\Pi$ is infinite then taking limits superior as $r \to \infty$ in the above shows than $\Pi$ must be of maximal type.

    If $\Delta_\Pi$ is finite then let $\Delta > \Delta_\Pi$. By definition of the limit superior there is an $r_0 > 0$ such that, for all $r \geq r_0$, we have
    \begin{equation*}
        n_\Pi(r) \leq \Delta r^{\nu_\Pi} = \Delta r^{\rho_\Pi}. 
    \end{equation*}
    As demonstrated in the proof of \Cref{thm:exponent-of-convergence-Weierstraß-product}, we can find a constant $C_{\lambda_\Pi} > 0$ such that for all $r > 0$ we have
    \begin{equation*}
        n_\Pi(r) \leq C_{\lambda_\Pi} \Delta r^{\rho_\Pi}
    \end{equation*}
    The genus $\mu$ of the canonical product satisfies $\mu < \lambda_\Pi < \mu + 1$, since $\lambda_\Pi$ is not an integer. Therefore, by \Cref{lem:canonical-product-estimate}, for all $r > 0$ and $\vert z \vert = r$ we have
    \begin{align*}
        \log \vert \Pi(z) \vert &< k_\mu r^\mu \left( \int_0^r \frac{n_\Pi(t)}{t^{\mu + 1}} \diff t + r \int_r^\infty \frac{n_\Pi(t)}{t^{\mu + 2}} \diff t \right) \leq \\
        &\leq k_\mu r^\mu C_{\lambda_\Pi} \Delta \left( \int_0^r t^{\rho_\Pi - \mu - 1} \diff t + r \int_r^\infty t^{\rho_\Pi - \mu - 2} \diff t \right) = \\
        &= k_\mu r^\mu C_{\lambda_\Pi} \Delta \left( \frac{r^{\rho_\Pi - \mu}}{\rho_\Pi - \mu} - \frac{r^{\rho_\Pi - \mu}}{\rho_\Pi - \mu - 1} \right) = k_\mu C_{\lambda_\Pi} \Delta \left( \frac{1}{\rho_\Pi - \mu} - \frac{1}{\rho_\Pi - \mu - 1} \right) r^{\rho_\Pi}.
    \end{align*}
    Dividing by $r^{\rho_\Pi}$, taking limits superior as $r \to \infty$ and letting $\Delta \searrow \Delta_\Pi$ therefore yields $\tau_\Pi \leq C \Delta_\Pi$ for a constant $C > 0$. Combining with \eqref{eq:density-smaller-constant-type}, we get
    \begin{equation*}
        C_1 \tau_\Pi \leq \Delta_\Pi \leq C_2 \tau_\Pi
    \end{equation*}
    for constants $C_1, C_2 > 0$. This shows that, if $\Delta_\Pi \in (0, \infty)$, then $\Pi$ is of normal type, and if $\Delta_\Pi = 0$, then $\Pi$ is of minimal type.
\end{proof}
\section{Hadamard's Theorem} \label{sec:hadamards-theorem}

Our main goal in this section will be to refine the factorization given by the Weierstraß Factorization Theorem for entire functions of finite order. Such a refinement is given by Hadamard's Theorem, the proof of which relies on the following lemma, which can be considered as a version of the maximum modulus principle applied to the real part of a holomorphic function.

The contents of this section closely follow the results found in Segal \cite{segal-complex-analysis}.

\begin{lemma}[Borel--Carathéodory] \label{lem:borel-caratheodory}
    Let $G \subseteq \C$ be a domain, $R > 0$ and suppose $\cl{B_R(0)} \subseteq G$ and $f \in H(G)$. Define $ A_f(r) \coloneqq \max_{\vert z \vert = r} \Re f(z) $, then, for $0 < r < R$,
    \begin{equation} \label{eq:borel-caratheodory-1}
        M_f(r) \leq \frac{2r}{R - r} A_f(R) + \frac{R + r}{R - r} \vert f(0) \vert
    \end{equation}
    and, if additionally $A_f(R) \geq 0$, then for any $n \in \N$
    \begin{equation} \label{eq:borel-caratheodory-2}
        M_{f^{(n)}}(r) \leq \frac{2^{n+2} n! R}{(R - r)^{n+1}} (A_f(R) + \vert f(0) \vert).
    \end{equation}
\end{lemma}

\begin{proof}
    If $f$ is constant, there is nothing to show.

    We assume $f$ being non-constant. First, for $r > 0$ we have
    \begin{align*}
        A_f(r) = \max_{\vert z \vert = r} \Re f(z) = \log \max_{\vert z \vert = r} \exp {\Re f(z)} = \log \max_{\vert z \vert = r} \vert \exp f(z) \vert = \log M_{\exp f}(r)
    \end{align*}
    and since $M_{\exp f}$ is strictly increasing and continuous therefore so is $A_f$.

    We will show (\ref{eq:borel-caratheodory-1}) by considering two cases. First assume $f$ is non-constant and $f(0) = 0$. Then by the above we have $A_f(R) > A_f(0) = 0$. Define
    \begin{equation*}
        \phi(z) \coloneqq \frac{f(z)}{2 A_f(R) - f(z)} \in H(B_R(0)).
    \end{equation*}
    Note that we have $\phi(0) = 0$ and
    \begin{align*}
        \vert \phi(z) \vert^2 &= \phi(z) \overline{\phi(z)} = \frac{\vert f(z) \vert^2}{(2 A_f(R) - f(z))(2 A_f(R) - \overline{f(z)})} = \\
        &= \frac{\vert f(z) \vert^2}{(2 A_f(R) - \Re f(z))^2 + \vert f(z) \vert^2 - (\Re f(z))^2} \leq 1,
    \end{align*}
    since clearly $2 A_f(R) - \Re f(z) \geq \Re f(z)$. Now, since $\phi(zR) \in H(\D)$ and $\vert \phi(zR) \vert \leq 1$ for $z \in \D$, Schwartz' Lemma implies $\vert \phi(zR) \vert \leq \vert z \vert$ for $z \in \D$. Therefore, for $z \in B_R(0)$ we have $\vert \phi (z) \vert \leq \frac{\vert z \vert}{R}$, and for $z \in B_r(0)$ we have $\vert \phi (z) \vert < \frac{r}{R} < 1$. Since
    \begin{equation*}
        f(z) = 2 A_f(R) \frac{\phi(z)}{1 + \phi(z)}
    \end{equation*}
    we obtain
    \begin{align*}
        \vert f(z) \vert &= 2 A_f(R) \left\vert \sum_{n=1}^\infty (-1)^n \phi(z)^n \right\vert \leq 2 A_f(R) \sum_{n=1}^\infty \left( \frac{r}{R} \right)^n = \frac{2 r}{R - r} A_f(R).
    \end{align*}

    Now assume that $f$ is non-constant and $f(0) \neq 0$. Set $g(z) \coloneqq f(z) - f(0)$, then by the above we have, for $w \in B_r(0)$,
    \begin{align*}
        \vert f(w) \vert - \vert f(0) \vert \leq \vert g(w) \vert \leq M_g(r) \leq \frac{2 r}{R - r} A_g(R) \leq \frac{2r}{R - r} A_f(R) - \frac{2r}{R - r} \Re f(0).
    \end{align*}
    Since $-\Re f(0) \leq \vert f(0) \vert$, this implies
    \begin{align*}
        \vert f(w) \vert \leq \frac{2r}{R - r}A_f(R) + \frac{2r}{R - r} \vert f(0) \vert + \vert f(0) \vert \leq \frac{2r}{R - r} A_f(R) + \frac{R + r}{R - r} \vert f(0) \vert,
    \end{align*}
    thus proving (\ref{eq:borel-caratheodory-1}).

    To show (\ref{eq:borel-caratheodory-2}), let $z \in \partial B_r(0)$ and set $s \coloneqq \frac{R - r}{2}$, then by Cauchy's integral formula
    \begin{equation*}
        f^{(n)}(z) = \frac{n!}{2 \pi i} \oint_{\partial B_s(z)} \frac{f(\zeta)}{(\zeta - z)^{n+1}} \diff \zeta.
    \end{equation*}
    Replacing $r$ with $\frac{R + r}{2} < R$ in (\ref{eq:borel-caratheodory-1}) we get
    \begin{align*}
        \vert f^{(n)}(z) \vert &\leq \frac{n!}{2 \pi} 2 \pi \frac{R - r}{2} \left( \frac{2}{R - r} \right)^{n+1} \left( \frac{2 (R + r)/2}{R - (R+r)/2} A_f(R) + \frac{R + (R + r)/2}{R - (R + r)/2} \vert f(0) \vert \right) \leq \\
        &\leq \frac{2^{n+1} n!}{(R - r)^{n+1}} \frac{R - r}{2} \left( 2 \frac{R+r}{R-r} A_f(R) + \frac{3R + r}{R - r} \vert f(0) \vert \right) \leq \\
        &\leq \frac{2^{n+1} n!}{(R - r)^{n+1}} \left( (R+r) A_f(R) + \frac{3R + r}{2} \vert f(0) \vert \right) \leq \frac{2^{n+2} n! R}{(R - r)^{n+1}} (A_f(R) + \vert f(0) \vert).
    \end{align*}
\end{proof}

\begin{theorem}[Hadamard]
    Let $f \in H(\C)$ be not identically zero, of finite order and let $m \in \N_0$ be the order of the zero of $f$ at the origin. Then there exists a polynomial $Q$ with $\deg Q \leq \rho_f$ such that
    \begin{equation}
        f(z) = z^m e^{Q(z)} \Pi(z),
    \end{equation}
    where $\Pi$ denotes the canonical product formed from the zeros of $f$.
\end{theorem}

Hadamard's Theorem can also be formulated more concisely: Let $f \in H(\C)$ be of finite order, then $\mu_f \leq \rho_f$, that is the genus of $f$ does not exceed its order.

\begin{proof}
    We may assume $f(0) \neq 0$. By the Weierstraß Factorization Theorem we have $f(z) = e^{Q(z)} \Pi(z)$ for some $Q \in H(\C)$ and a Weierstraß product $\Pi$. Since
    $$ \mu_\Pi \leq \lambda_\Pi = \lambda_f \leq \rho_f $$
    we may take $\Pi$ to be a canonical product. Set $\nu \coloneqq \lfloor \rho_f \rfloor$, then taking logarithms in the above factorization and differentiating $\nu + 1$ times we get
    \begin{align*}
        \deriv{z}{\nu} \left( \frac{f'(z)}{f(z)} \right)
        &= Q^{(\nu+1)}(z) + \deriv{z}{\nu + 1} \log \prod_{m \in M} W \left( \frac{z}{z_m}; \mu_\Pi \right) = \\
        &= Q^{(\nu+1)}(z) + \deriv{z}{\nu + 1} \sum_{m \in M} \left( \log \left( 1 - \frac{z}{z_m} \right) + \sum_{k=1}^{\mu_\Pi} \frac{z^k}{k} \right) = \\
        &= Q^{(\nu+1)}(z) - \sum_{m \in M} \frac{1}{z_m} \deriv{z}{\nu} \left( 1 - \frac{z}{z_m} \right)^{-1} = \\
        &= Q^{(\nu+1)}(z) - \nu! \sum_{m \in M} \frac{1}{(z_m - z)^{\nu + 1}}.
    \end{align*}
    We now aim to show that $Q^{(\nu + 1)}$ is identically zero. For $R > 0$ we set
    $$ g_R(z) \coloneqq \frac{f(z)}{f(0)} \prod_{\substack{m \in M \\ \vert z_m \vert \leq R}} \left(1 - \frac{z}{z_m} \right)^{-1} \in H(\C). $$
    Note that $g_R(0) = 1$. For $\vert z \vert = 2R$ and $\vert z_m \vert \leq R$ we have $\vert 1 - z / z_m \vert \geq \vert z \vert / \vert z_m \vert - 1 \geq 2R/R - 1 = 1$, therefore
    $$ M_{g_R}(2R) \leq \max_{\vert z \vert = 2R} \left\vert \frac{f(z)}{f(0)} \right\vert = \O(\exp (2R)^\alpha), \quad \textrm{as } R \to \infty $$
    for all $\alpha > \rho_f$. By the maximum modulus theorem this also holds for $\vert z \vert \leq 2R$, and in particular we have
    $$ \log M_{g_R}(R) = \O(R^\alpha). $$
    Note that $g_R$ has no zeros in $\cl{B_R(0)}$, therefore we can define $h_R(z) \coloneqq \log g_R(z)$, choosing the branch of the logarithm for which $h_R(0) = 0$, then $h_R$ is holomorphic on some domain containing $\cl{B_R(0)}$. We have
    $$ A_{h_R}(R) = \max_{\vert z \vert = R} \Re h_R(z) = \max_{\vert z \vert = R} \log \vert g_R(z) \vert = \log M_{g_R}(R) = \O(R^\alpha). $$
    Since by the maximum modulus theorem $\vert g_R(z) \vert \geq \vert g_R(0) \vert = 1$ we have $\Re h_R(z) \geq 0$. Invoking the Borel-Carathéodory lemma we therefore get, for $0 < r < R$ and sufficiently large $R$,
    $$ M_{h_R^{(\nu + 1)}}(R) \leq \frac{2^{\nu + 3} (\nu + 1)! R}{(R - r)^{\nu + 2}} A_{h_R}(R) \leq \frac{K_0 2^{\nu + 3} (\nu + 1)! R^{\alpha + 1}}{(R - r)^{\nu + 2}}, $$
    where $K_0 > 0$ is some constant depending only on $\alpha$. With $r \coloneqq R / 2$ it follows that
    $$ M_{h_R^{(\nu + 1)}}(R) \leq K_1 R^{\alpha - (\nu + 1)}, $$
    where $K_1 > 0$ is some constant depending only on $\alpha$ and $\nu$. Note that we have
    $$ h_R^{(\nu + 1)}(z) = \deriv{z}{\nu} \left( \frac{f'(z)}{f(z)} \right) + \nu! \sum_{\substack{m \in M \\ \vert z_m \vert \leq R}} \frac{1}{(z_m - z)^{\nu + 1}} = Q^{(\nu + 1)}(z) - \nu! \sum_{\substack{m \in M \\ \vert z_m \vert > R}} \frac{1}{(z_m - z)^{\nu + 1}}. $$
    For $\vert z \vert = R / 2$ and $z_m > R$ we have $\vert z_m - z \vert \geq \vert z_m \vert - R / 2 \geq \vert z_m \vert / 2$, therefore
    $$ \vert Q^{(\nu+1)}(z) \vert \leq \nu! \sum_{\substack{m \in M \\ \vert z_m \vert > R}} \frac{1}{\vert z_m - z \vert^{\nu + 1}} + \vert h_R^{(\nu + 1)}(z) \vert \leq K_2 \sum_{\substack{m \in M \\ \vert z_m \vert > R}} \frac{1}{\vert z_m \vert^{\nu + 1}} + K_1 R^{\alpha - (\nu + 1)}, $$
    where $K_2$ is some constant depending only on $\nu$. By the maximum modulus principle this holds for all $\vert z \vert \leq R / 2$. Since $\nu + 1 > \rho_f \geq \lambda_f$ both terms on the right side tend to zero as $R \to \infty$, if $\rho_f < \alpha < \nu + 1$. Therefore $Q^{(\nu + 1)}$ vanishes identically, thus $Q$ is a polynomial of degree at most $\nu \leq \rho_f$.
\end{proof}

\begin{remark} \label{rem:consequences-hadamard}
    An immediate consequence of Hadamard's Theorem is that entire functions of order $\rho < 1$ are of genus $0$, and are therefore -- like polynomials -- solely determined by their roots and some scaling factor. Indeed, in this case the polynomial $Q$ in the factorization is constant and the canonical product $\Pi$ is of genus $0$ as well. In particular, any non-vanishing function of order zero is constant.

    Furthermore, we can deduce that non-polynomial functions of order zero have no Picard exceptional value: Let $f$ denote such a function and assume towards a contradiction that $f$ assumes some value $a \in \C$ only finitely often. Then, by Hadamard's Theorem, we can write
    $$ f(z) - a = C z^m \Pi(z) $$
    for some constant $C > 0, m \in \N_0$ and a canonical product $\Pi$. Since $a$ is assumed only finitely often, the canonical product $\Pi$ must be finite and therefore $f$ a polynomial, a contradiction.
    
    As shown in \Cref{rem:order-zero} the entire function
    $$ f(z) \coloneqq \sum_{k=0}^\infty \frac{z^k}{(k^2)!} $$
    is an example of a non-polynomial, entire function of order zero and must therefore, by the above reasoning, assume all values in $\C$ infinitely often.
\end{remark}

Hadamard's Theorem has a few consequences for entire functions of finite order, which we will explore further. Combining it with \Cref{thm:exponent-of-convergence-Weierstraß-product} allows us to prove two results regarding functions of finite, non-integer order.

\begin{theorem} \label{thm:finite-non-integer-order-equals-exponent-of-convergence}
    Let $f \in H(\C)$ be of finite, non-integer order. Then $\rho_f = \lambda_f$.
\end{theorem}

\begin{proof}
    By \Cref{thm:inequality-order-exponent-of-convergence} we have $\lambda_f \leq \rho_f$. Invoking Hadamard's Theorem we can write
    $$ f(z) = z^m e^{Q(z)} \Pi(z) $$
    for some $m \in \N_0$, a polynomial $Q$ with $\deg Q \leq \rho_f$ and a canonical product $\Pi$. Since $\rho_f$ is not an integer we have $\deg Q \leq \lfloor \rho_f \rfloor < \rho_f$. By \Cref{prop:order-exponential-polynomial} $e^Q$ has order at most $\deg Q$ and by \Cref{thm:exponent-of-convergence-Weierstraß-product} $\Pi$ has order $\lambda_f$, therefore using \Cref{prop:order-sum-product-estimate} we obtain
    $$ \rho_f \leq \max \{ 0, \deg Q, \lambda_f \} = \lambda_f \leq \rho_f, $$
    since $\rho_f \leq \max \{ \deg Q, \lambda_f \} = \deg Q < \rho_f$ would be a contradiction, and we get $\rho_f = \lambda_f$.
\end{proof}

\begin{theorem} \label{thm:finite-non-integer-order-infinite-zeros}
    Let $f \in H(\C)$ be of finite, non-integer order. Then $f$ has infinitely many zeros.
\end{theorem}

\begin{proof}
    By \Cref{thm:finite-non-integer-order-equals-exponent-of-convergence} we have $\rho_f = \lambda_f$. Since $\rho_f$ is not an integer it follows that $\lambda_f > 0$, which implies that $f$ has infinitely many zeros.
\end{proof}

\begin{theorem}[Borel] \label{thm:existence-borel-exceptional-values}
    Let $f \in H(\C)$ be of integer order $\rho_f > 0$. Then for any $a \in \C$ we have $\lambda_f^{(a)} = \rho_f$, except possibly for one value of $a$.
\end{theorem}

\begin{proof}
    Note that for any $w \in \C$ we have
    $$ \lambda_f^{(w)} = \lambda_{f-w} \leq \rho_{f-w} = \rho_f. $$
    Assume towards a contradiction that there exist $a, b \in \C$ such that $\lambda_f^{(a)} < \rho_f$ and $\lambda_f^{(b)} < \rho_f$. By Hadamard's Theorem we can write
    $$ \alpha(z) \coloneqq f(z) - a = z^{m_1} e^{Q_1(z)} \Pi_1(z), \quad \beta(z) \coloneqq f(z) - b = z^{m_2} e^{Q_2(z)} \Pi_2(z), $$
    where $m_1, m_2 \in \N_0$, $Q_1, Q_2$ are polynomials of degree at most $\rho_f$ and $\Pi_1, \Pi_2$ denote appropriate canonical products. We have
    $$ \deg Q_1 \leq \rho_f = \rho_\alpha \leq \max \{ 0, \deg Q_1, \rho_{\Pi_1} \}, $$
    and since by \Cref{thm:exponent-of-convergence-Weierstraß-product} $\rho_{\Pi_1} = \lambda_f^{(a)} < \rho_f$, it follows that $\deg Q_1 = \rho_f$. The analogous argument for $Q_2$ yields $\deg Q_2 = \rho_f$. Subtracting $\beta$ from $\alpha$ yields
    \begin{equation} \label{thm:existence-borel-exceptional-values:difference-of-functions}
    \begin{aligned}
                        && b - a &= z^{m_1} e^{Q_1(z)} \Pi_1(z) - z^{m_2} e^{Q_2(z)} \Pi_2(z) \\
        \Leftrightarrow && (b - a) e^{-Q_2(z)} &= z^{m_1} e^{Q_1(z) - Q_2(z)} \Pi_1(z) - z^{m_2} \Pi_2(z) \\
        \Leftrightarrow && (b - a) e^{-Q_2(z)} + z^{m_2} \Pi_2(z) &= z^{m_1} e^{Q_1(z) - Q_2(z)} \Pi_1(z)
    \end{aligned}
    \end{equation}
    Since $\rho_{\Pi_2} = \lambda_f^{(b)} < \rho_f$ and $-Q_2$ is of degree $\rho_f$, the left-hand side is of order $\rho_f$ and thus so is the right-hand side. Now similarly, since $\rho_{\Pi_1} = \lambda_f^{(b)} < \rho_f$ and the right-hand side has order $\rho_f$, we have that $Q_1 - Q_2$ is of degree $\rho_f$.
    
    Differentiating the first equation in \eqref{thm:existence-borel-exceptional-values:difference-of-functions} gives
    \begin{equation*}
    \begin{aligned}
        && 0 &= m_1 z^{m_1 - 1} e^{Q_1(z)} \Pi_1(z) + z^{m_1} Q_1'(z) e^{Q_1(z)} \Pi_1(z) + z^{m_1} e^{Q_1}(z) \Pi_1'(z) \\
        && &\quad\quad - m_2 z^{m_2 - 1} e^{Q_2(z)} \Pi_2(z) - z^{m_2} Q_2'(z) e^{Q_2(z)} \Pi_2(z) - z^{m_2} e^{Q_2(z)} \Pi_2'(z) \\
        \Leftrightarrow && 0 &= e^{Q_1(z)} (m_1 z^{m_1 - 1} \Pi_1(z) + z^{m_1} Q_1'(z) \Pi_1(z) + z^{m_1} \Pi_1'(z)) \\
        && &\quad\quad - e^{Q_2(z)} (m_2 z^{m_2 - 1} \Pi_2(z) + z^{m_2} Q_2'(z) \Pi_2(z) + z^{m_2} \Pi_2'(z))
    \end{aligned}
    \end{equation*}
    By \Cref{prop:order-derivative} we have $\rho_{\Pi_1'} = \rho_{\Pi_1}$ and $\rho_{\Pi_2'} = \rho_{\Pi_2}$, therefore the coefficients of $e^{Q_1}$ and $e^{Q_2}$ have order less than $\rho_f$. By Hadamard's Theorem we therefore get
    \begin{equation*}
        e^{Q_1(z)} z^{m_3} e^{Q_3(z)} \Pi_3(z) = e^{Q_2(z)} z^{m_4} e^{Q_4(z)} \Pi_4(z)
    \end{equation*}
    where $m_3, m_4 \in \N_0, Q_3, Q_4$ are polynomials of degree less than or equal to $\rho_f - 1$ (since $\rho_f$ is an integer) and $P_3, P_4$ denote canonical products. Rewriting the above we get
    \begin{equation*}
        z^{m_3} e^{Q_1(z) + Q_3(z)} \Pi_3(z) = z^{m_4} e^{Q_2(z) + Q_4(z)} \Pi_4(z)
    \end{equation*}
    Since both sides are equal, in particular they must share the same zeros and multiplicities, therefore we have $m_3 = m_4$ and $\Pi_3 = \Pi_4$. But this implies
    \begin{equation*}
    \begin{aligned}
        && Q_1(z) + Q_3(z) &= Q_2(z) + Q_4(z) \\
        \Leftrightarrow && Q_2(z) - Q_1(z) &= Q_3(z) - Q_4(z)
    \end{aligned}
    \end{equation*}
    and by the above the left side is of degree $\rho_f$, whereas the right side is of degree less than $\rho_f$, a contradiction.
\end{proof}

In the context of Borel's Theorem, we refer to such an exceptional value as a \emph{Borel exceptional value}. The theorem thus shows that entire functions of integer order have at most one Borel exceptional value.

Borel's Theorem together with \Cref{thm:finite-non-integer-order-infinite-zeros} allow us to obtain a stronger version of Picard's Little Theorem for functions of finite, positive order:

\begin{corollary}
    Let $f \in H(\C)$ with $\rho_f \in (0, \infty)$. Then $f$ assumes at most one value only finitely often.
\end{corollary}

\begin{proof}
    If $\rho_f$ is an integer then and $f$ takes on some value $a \in \C$ only on finitely often, then $\lambda_f^{(a)} = 0 < \rho_f$ and by Borel's Theorem there is at most one value of $a$ for which this holds.

    If $\rho_f$ is not an integer, then for any $w \in \C$ the function $f(z) - w$ is also of finite, non-integer order and therefore has infinitely many zeros by \Cref{thm:finite-non-integer-order-infinite-zeros}. Thus $f$ assumes all values infinitely often.
\end{proof}

Hadamard's Theorem can also be used to obtain a generalization of \Cref{thm:canonical-product-type-density} which holds for general entire functions of non-integer order, not just canonical products. This requires the following result, the proof of which relies on finding a lower bound on the modulus of arbitrary entire functions and can be found in Levin \cite{levin-distribution-of-zeros} as Theorem 12(a).

\begin{proposition} \label{prop:order-type-product-equality}
    Let $f, g \in H(\C)$ be of finite order. If $\rho_f < \rho_g$, then it holds that $ \rho_{fg} = \rho_g $ and\footnote{This holds in the sense that if $g$ is of maximal or minimal type, then $fg$ is of maximal or minimal type respectively, and, if $g$ is of normal type, then $\tau_{fg} = \tau_g$.} $\tau_{fg} = \tau_g$.
\end{proposition}

\begin{theorem} \label{thm:entire-function-type-density}
    Let $f \in H(\C)$ be of finite, non-integer order. Then $f$ is of maximal, minimal or normal type according to whether $\Delta_f$ is equal to infinity, zero, or a positive, real number.
\end{theorem}

\begin{proof}
    By Hadamard's Theorem we can write
    $$ f(z) = z^m e^{Q(z)} \Pi(z). $$
    The first factor is of order zero, the second is at most of order $\lfloor \rho_f \rfloor < \rho_f$ and the last is of order $\rho_f$. Therefore, by \Cref{prop:order-type-product-equality} we have  $\tau_f = \tau_\Pi$. Since the order of $\Pi$ is not an integer and clearly $\Delta_f = \Delta_\Pi$, the assertion now follows from \Cref{thm:canonical-product-type-density}.
\end{proof}

In general, if $f \in H(\C)$ is of finite order with representation
$$ f(z) = e^{Q(z)} \Pi(z), $$
where $Q$ is a polynomial and $\Pi$ is a canonical product, the order and type of $f$ may be determined by either of the factors. However, \Cref{thm:finite-non-integer-order-equals-exponent-of-convergence,thm:entire-function-type-density} show that if $\rho_f$ is not an integer, then the growth of $f$ is determined solely by the canonical product $\Pi$, and therefore by the zeros of $f$.
\chapter{Composition of Entire Functions}
\label{ch:composition-entire-functions}

As seen by \Cref{prop:algebraic-properties-order}, the sum or product of two entire functions of finite order is again of finite order. This is not the case when composition is involved; indeed, consider the function $f(z) \coloneqq \exp \exp z \in H(\C)$, which is of infinite order, yet is the composition of two functions of order $1$. 

\section{Pólya's Theorem}
\label{sec:polyas-theorem}

Necessary conditions for the order of a composition to be finite will be given by Pólya's Theorem, the proof of which relies on a result that was first proven by Harald Bohr. While Bloch's Theorem dealt with disks contained in the image of holomorphic functions, we now focus on circles contained in such images.

\begin{proposition} \label{prop:existence-circle-maximal-radius}
    Let $G \subset \C$ be a bounded domain with $0 \in G$. Then the set
    \begin{equation} \label{eq:set-circle-maximal-radius}
        S \coloneqq \{ r \geq 0 : \partial B_r(0) \subseteq \cl{G} \}
    \end{equation}
    has a positive maximum; that is, $\cl{G}$ contains a circle of positive, maximal radius\footnote{The circle of radius $0$ is understood as the singleton $\{ 0 \}$.}.
\end{proposition}

\begin{proof}
    Since $G$ is open and contains $0$ there is some $t > 0$ such that $B_t(0) \subseteq G$. Since $\partial B_t(0) \subset \cl{B_t(0)} \subseteq \cl{G}$ we thus have $t \in S$ and $S$ is non-empty. It now suffices to show that $S$ is compact, since then it would contain its maximum $m \geq t$. Clearly $S$ is bounded from below by $0$ and since $\cl{G}$ is bounded it must also be bounded from above. Thus it remains to show that $S$ is closed.

    Let $(r_n)_{n \in \N}$ be a convergent sequence in $S$ with limit $r \in [0, \infty)$. If $r = 0$, then $r \in S$. Otherwise we claim that $\partial B_r(0) \subseteq \cl{G}$. Let $z \in \partial B_r(0)$, and set $z_n \coloneqq r_n z / r$ for $n \in \N$. Now $\vert z_n \vert = r_n$, thus $z_n \in \partial B_{r_n}(0) \subseteq \cl{G}$ and therefore $(z_n)_{n \in \N}$ is a sequence in $\overline{G}$. Clearly $z_n \to z$ and since $\cl{G}$ is closed we therefore have $z \in \overline{G}$. Since $z$ was arbitrary, this concludes the claim, we have $r \in S$ and thus $S$ is closed.
\end{proof}

If $G \subset \C$ is a bounded domain and $f \in H(G) \cap C(\cl{G}; \C)$ is non-constant with $f(0) = 0$, then $f(G)$ is a bounded domain with $0 \in f(G)$. This, together with \Cref{prop:existence-circle-maximal-radius}, justifies the following definition:

\begin{definition}
    Let $f \in H(\D) \cap C(\cl{\D}; \C)$ be non-constant with $f(0) = 0$, then we define $r_f$ as the positive maximum of the set $S$ in \eqref{eq:set-circle-maximal-radius}, where $G = f(\D)$.
\end{definition}

Bohr's Theorem now asserts that $r_f$ does (almost) not depend on $f$ itself:

\begin{theorem}[Bohr] \label{thm:bohr}
    There exists a function ${\phi : (0, 1) \to (0, \infty)}$ such that for any $\theta \in (0, 1)$ and $f \in H(\D) \cap C(\cl{\D}; \C)$ with $f(0) = 0$ and $M_f(\theta) = 1$ it holds that $r_f \geq \phi(\theta)$.
\end{theorem}

\begin{proof}
    Suppose $f$ satisfies the hypothesis of the theorem, let $\varepsilon > 0$ and set $R_\varepsilon \coloneqq r_f + \varepsilon$. By definition of $r_f$, for all $r \geq R_\varepsilon$ there exists some point $w_r \in \partial B_r(0)$ with $w_r \notin f(\cl{\D})$. Choose such points $w_{R_\varepsilon}, w_{2 R_\varepsilon}$ and define
    $$ h(z) \coloneqq \frac{f(z) - w_{R_\varepsilon}}{w_{2 R_\varepsilon} - w_{R_\varepsilon}} \in H(\D; \C \setminus \{ 0, 1 \}). $$
    Since
    $$ \vert h(0) \vert = \left\vert \frac{f(0) - w_{R_\varepsilon}}{w_{2 R_\varepsilon} - w_{R_\varepsilon}} \right\vert \leq \frac{\vert w_{R_\varepsilon} \vert}{\vert \vert w_{2 R_\varepsilon} \vert - \vert w_{R_\varepsilon} \vert \vert} = \frac{R_\varepsilon}{2 R_\varepsilon - R_\varepsilon} = 1, $$
    by Schottky's Theorem we have $\vert h(z) \vert \leq \psi(\theta, 1)$ for all $\vert z \vert \leq \theta$. Therefore
    $$ \vert f(z) \vert - R_\varepsilon \leq \vert g(z) - w_{R_\varepsilon} \vert \leq \vert w_{2 R_\varepsilon} - w_{R_\varepsilon} \vert \psi(\theta, 1) \leq 3 R_\varepsilon \psi(\theta, 1) $$
    and thus $\vert f(z) \vert \leq R_\varepsilon + 3 R_\varepsilon \psi(\theta, 1)$ for all $\vert z \vert \leq \theta$. Using the hypothesis that $M_f(\theta) = 1$ and the maximum modulus principle we obtain $1 \leq R_\varepsilon + 3 R_\varepsilon \psi(\theta, 1)$, and letting $\varepsilon \to 0$ we have $1 \leq r_f + 3 r_f \psi(\theta, 1)$. Thus
    $$ r_f \geq \frac{1}{1 + 3 \psi(\theta, 1)}$$
    and defining $\phi(\theta)$ as the right-hand side establishes the assertion.
\end{proof}

\begin{theorem}[Pólya] \label{thm:polya}
    Let $g, h \in H(\C)$ be non-constant. For the order of $g \circ h$ to be finite, it must hold that either
    \begin{enumerate}[i.]
        \item $h$ is a polynomial and $\rho_g < \infty$, or
        \item $h$ is not a polynomial, $\rho_h < \infty$ and $\rho_g = 0$.
    \end{enumerate}
\end{theorem}

\begin{proof}
    Without loss of generality we can assume $h(0) = 0$; otherwise we just consider $h_0(z) \coloneqq h(z) - h(0)$ and $g_0(w) \coloneqq g(w + h(0))$. Set $f \coloneqq g \circ h$ and define
    \begin{equation*}
        k_r(z) \coloneqq \frac{h(rz)}{M_h(r / 2)} \in H(\D) \cap C(\cl{\D}; \C), \quad \textrm{for } r > 0.
    \end{equation*}
    Note that by definition we have $M_{k_r}(1/2) = 1$ and $k_r(0) = 0$, thus by Bohr's Theorem there is some constant $C > 0$ and an $R > C M_h(r/2)$ such that $\partial B_{R / M_h(r/2)}(0) \subseteq k_r(\cl{\D})$ and thus $\partial B_R(0) \subseteq h(\cl{B_r(0)})$. By the maximum modulus principle, $g$ assumes its maximum modulus over $\cl{B_R(0)}$ at some $w_0 \in \partial B_R(0)$. By the above there is a $z_0 \in \cl{B_r(0)}$ with $h(z_0) = w_0$. Thus we get
    \begin{equation*}
        M_g(C M_h(r / 2)) < M_g(R) = \vert g(w_0) \vert = \vert g(h(z_0)) \vert = \vert f(z_0) \vert \leq M_f(r).
    \end{equation*}
    Assuming $\rho_f < \infty$, we have $M_f(r) < K \exp(r^\alpha)$ for every $\alpha > \rho_f$, with suitable $K > 0$. Consider the power series expansion $h(z) = \sum_{n=0}^\infty a_n z^n$. Let $a_m$ denote any non-zero coefficient; note that since $h(0) = 0$ we have $m \geq 1$. By Cauchy's integral formula we have, for all $s > 0$,
    \begin{equation*} \label{eq:proof-polya-cauchy-inequality}
        \vert a_m \vert = \left\vert \frac{h^{(m)}(0)}{m!} \right\vert = \frac{1}{2\pi} \left\vert \oint_{\partial B_{s}(0)} \frac{h(\zeta)}{\zeta^{n+1}} \diff \zeta \right\vert \leq \frac{M_h(s)}{s^m} \tag{\textasteriskcentered}
    \end{equation*}
    and thus
    \begin{equation*}
        M_g(C \vert a_m \vert (r/2)^m) \leq M_g(C M_h(r/2)) < M_f(r) < K \exp(r^\alpha), \quad \textrm{for all } r > 0.
    \end{equation*}
    Replacing $(r/2)^m$ with $r$ we obtain $\rho_g \leq \alpha / m $. If $h$ is not a polynomial we may let $m \to \infty$, thus $\rho_g = 0$.

    Now consider $g(z) = \sum_{n=0}^\infty b_n z^n$. Replacing $h$ with $g$ in (\ref{eq:proof-polya-cauchy-inequality}) we obtain $\vert b_n \vert s^n \leq M_g(s)$ for all $s > 0$ and $n \geq 1$ and thus
    \begin{equation*}
        \vert b_n \vert (C M_h(r/2))^n \leq M_g(C M_h(r/2)) < M_f(r) < K \exp(r^\alpha),
    \end{equation*}
    which implies $\rho_h \leq \alpha < \infty$.
\end{proof}

\begin{theorem}[Thron] \label{thm:thron}
    Let $g \in H(\C)$ be transcendental, $\rho_g < \infty$ and suppose that $g$ assumes some value $w \in \C$ only finitely often. Then there exists no $f \in H(\C)$ with $f \circ f = g$.
\end{theorem}

\begin{proof}
    Assume towards a contradiction that $f \in H(\C)$ with $f \circ f = g$ exists. Since $g$ is not a polynomial, $f$ is not a polynomial either. Thus Pólya's Theorem implies $\rho_f = 0$.
    
    Consider the sets
    $$ Z \coloneqq f^{-1}(\{ w \}), \quad Z' \coloneqq \bigcup_{z \in Z} f^{-1}(\{ z \}). $$
    By definition, for each $z' \in Z'$ there is some $z \in Z$ with $z' \in f^{-1}(\{ z \})$, thus
    $$ g(z') = f(f(z')) = f(z) = w. $$
    Our hypothesis on $g$ implies that $Z'$ must be finite. Since pre-images of singletons are disjoint, $Z'$ is a disjoint union, therefore
    $$ \sum_{z \in Z} \vert f^{-1}(\{ z \}) \vert = \left\vert \bigcup_{z \in Z} f^{-1}(\{ z \}) \right\vert = \vert Z' \vert < \infty $$
    and thus all points in $Z$ are only assumed finitely often by $f$. But by \Cref{cor:transcendental-every-value-inf}, $f$ assumes at most one value only finitely often; therefore $\vert Z \vert \leq 1$.

    If $Z = \emptyset$, then $h(z) \coloneqq f(z) - w$ is entire, of order $0$ and nowhere $0$. As discussed in \Cref{rem:consequences-hadamard}, as a consequence of Hadamard's Theorem we have that $h$ is constant and consequently so is $f$, a contradiction.

    If $Z = \{ z_0 \}$, then $h(z) \coloneqq f(z) - w$ has a single zero of finite order $n \in \N$ at $z_0$. Therefore we can write $h(z) = (z - z_0)^n p(z)$, where $p$ is entire, of order $0$ and nowhere $0$. Again, this implies that $p$ is constant, and therefore $f$ a polynomial, a contradiction.
\end{proof}

\begin{example}
    A natural application of Thron's Theorem is to set $g(z) \coloneqq \exp z$, which never assumes zero as a value. Indeed, this implies that there is $f \in H(\C)$ satisfying
    $$ f(f(z)) = \exp z. $$
    On the other hand, there does exist a real-analytic function satisfying the above. The construction of such a function is difficult, but was demonstrated by H. Kneser \cite{kneser-real-analytic-solution}.
\end{example}

\bibliography{literature.bib}
\bibliographystyle{abbrv}

\end{document}
